\documentclass[letterpaper,twocolumn,10pt]{article} %
\usepackage{amsmath}
\usepackage{usenix}
\newif\ifAnon\Anontrue
\newif\ifDraft\Draftfalse

\usepackage{todonotes}
\usepackage{cispa_defaults}
\usepackage{svg}
\usepackage{cpu_names}
\usepackage{stfloats}


\usepackage{subcaption}
\usepackage{multirow}
\usepackage{graphicx}
\usepackage{tcolorbox}
\usepackage{fontawesome5}
\usetikzlibrary{arrows.meta, positioning}
\usetikzlibrary{pgfplots.fillbetween} %
\usetikzlibrary{chains}



\newcommand{\code}[1]{\texttt{#1}\xspace}
\newcommand{\missingnumNum}{XXX\xspace}
\newcommand{\missingnumText}{\textcolor{red}{XXX}\xspace}
\newcommand{\paragrabf}[1]{\noindent \textbf{#1.}\ }
\newcommand{\td}[1]{\textcolor{orange}{\\TODO: #1}}

\newcommand{\AttackName}{PhotoFinish\xspace}


\newcommand{\PlatformCount}{38\xspace}
\newcommand{\IsaCount}{4\xspace}

\newcommand{\tprobe}{{\fontfamily{cmtt}\selectfont T0}\xspace}
\newcommand{\tref}{{\fontfamily{cmtt}\selectfont T1}\xspace}
\newcommand{\corS}{{\fontfamily{cmtt}\selectfont S}\xspace}
\newcommand{\corM}{{\fontfamily{cmtt}\selectfont M}\xspace}
\newcommand{\corO}{{\fontfamily{cmtt}\selectfont O}\xspace}
\newcommand{\rdelta}[1]{$\Delta_{#1}$}

\newlength{\maxheight}
\setlength{\maxheight}{\heightof{W}}
\newcommand{\baselinealign}[1]{%
\centering
\raisebox{0pt}[\maxheight][0pt]{#1}%
}

\newcommand{\takeaway}[1]{%
\begin{tcolorbox}[boxsep=1pt,left=2pt,right=2pt,top=1pt,bottom=1pt]
    \faLightbulb[regular]~\textbf{Takeaway:} #1
\end{tcolorbox}
}

\begin{document}

\title{\AttackName: Portable Timerless Probing of Microarchitectural Latencies}

\author{
  {\rm Tristan Hornetz, Eduard Ebert, Lorenz Hetterich, Lukas Gerlach, Michael Schwarz}\\
  CISPA Helmholtz Center for Information Security
}

\maketitle

\begin{abstract}
Timing measurements are a fundamental building block of microarchitectural attacks, enabling software to convert otherwise invisible latency differences into architectural observations.
Thus, restricting high-resolution timers is a common defense, especially in constrained environments such as web browsers.
Timerless probing primitives provide an alternative approach.
Instead of reading or constructing a timer, they observe architectural events that correlate with the microarchitectural state.
However, such primitives typically rely on platform-specific behavior, and none generalizes across the wider computing landscape.

In this paper, we present \AttackName, a portable timerless probing primitive that converts latency differences directly into an architectural value.
\AttackName races an operation under test against a fixed reference in two parallel threads that then write unique values to the same shared location, resulting in a latency-dependent architectural value.
\AttackName only requires parallel execution and shared memory, without constructing or reading a clock.
We evaluate \AttackName on \PlatformCount CPUs spanning four ISAs and demonstrate reliable cache-state classification across widely different microarchitectures, with comparable fidelity when compiled to WebAssembly and executed inside a browser.
We further use \AttackName in AES T-table and \EvictReload attacks, show that it applies beyond caches to TLB and execution-engine latencies, and demonstrate sufficient precision to identify the LLC slice of an address.
Our results demonstrate shared-memory races as a generic and precise replacement for explicit timing, showing that restricting timers alone does not prevent software from observing fine-grained microarchitectural latency differences.
\end{abstract}

\section{Introduction}
Software-based microarchitectural attacks remain a practical threat to the confidentiality and integrity of modern systems.
They exploit effects below the architectural abstraction, such as cache state~\cite{Osvik2006,Yarom2014Flush}, branch-predictor state~\cite{Evtyushkin2018,Lee2017Inferring}, TLB state~\cite{Gras2018,Lipp2022Prefetch}, execution-unit contention~\cite{aldaya2019port,Rokicki2022port}, or DRAM row-buffer state~\cite{Pessl2016}.
Transient-execution attacks additionally exploit instructions that execute transiently and leave observable microarchitectural traces~\cite{Kocher2019,Lipp2018meltdown,Canella2019A}.
Although these effects are not directly visible in a program's architectural state, they affect execution time.
Thus, attackers can infer information across security boundaries by observing small latency differences in otherwise ordinary software execution.

A central building block of many such attacks is a timing measurement primitive.
The attacker executes an operation whose latency depends on microarchitectural state and then classifies it as fast or slow.
In cache attacks, this operation is commonly a memory access~\cite{Yarom2014Flush,Gruss2016Flush,Rauscher2025Systematic}.
In other attacks, it may be a branch~\cite{Aciicmez2007predicting,Lee2017Inferring,Evtyushkin2018}, an instruction sequence contending for a shared execution resource~\cite{aldaya2019port,Hornetz2025PortPrint}, a page-table walk~\cite{Gruss2016Prefetch,Lipp2022Prefetch}, or an instruction's execution time~\cite{Andrysco2015subnormal}.
In all cases, the attacker needs a way to turn a microarchitectural latency difference into an observation that software can use.

Traditionally, attackers obtain this observation using an incrementing timer.
The attacker timestamps the operation before and after execution and compares the difference against a fixed threshold.
Consequently, a common defense strategy is to restrict access to high-resolution timers: browsers reduce timer resolution and introduce jitter~\cite{Schwarz2017Timers,VanGoethem2015,Shusterman2021prime}, or disable them in specific settings~\cite{Intel2021sgxdeveloper,Lipp2016,thomas2026riscy}.
These restrictions do not remove the underlying latency difference.
They only remove the most direct way to observe it.
Attackers can compensate by amplifying the targeted difference~\cite{Allan2016,Roettger2021SpectreJS,Schwarzl2021dynamic,Purnal2023showtime} or constructing software clocks such as counter threads~\cite{Schwarz2017Timers,Lipp2016,Gras2017aslr}.

An alternative is to avoid constructing a timer altogether.
Existing timerless primitives expose microarchitectural state through specialized instructions~\cite{Disselkoen2017prime,kim2022dprime,Yu2023Synchronization,Zhang2023MWAIT,Thomas2025Exfilstate,Hornetz2026tdxray} or environment-specific properties of the cache hierarchy~\cite{Guanciale2016Cache,Jin2025Ghostcache,Thomas2025Exfilstate,thomas2026riscy}.
These techniques demonstrate that explicit timestamps are not essential, but they replace the dependence on a timer with assumptions about a particular ISA, processor feature, or cache implementation.
This leaves a portability gap: existing timerless primitives do not provide a generic way to observe latency differences across architectures and restricted execution environments, such as browsers.

In this paper, we ask the following research question:

\emph{Can we build a generic primitive that observes microarchitectural latency differences without reading a timer and without relying on specialized architecture-specific features?}

To answer this question, we present \AttackName, a timerless probing primitive applicable to all major computer architectures and even to restricted environments like web browsers.
We base this primitive on an observation by Purnal et al.~\cite{Purnal2023showtime}: A race condition between two parallel threads can reveal microarchitectural timing differences when both threads write to a shared variable after finishing a workload.
Specifically, one thread, which we call the \emph{probing thread} \tprobe, performs the target operation, whose latency depends on the microarchitectural state.
In parallel, a \emph{reference thread} \tref performs a reference operation with known latency.
Finally, both threads write distinct values to the same shared result variable.
When both threads complete, the final value of this variable reveals which store serialized last, showing whether the target operation was faster or slower than the reference.
The key property of this approach is that the attacker does not need a timestamp.
The primitive only requires two threads, shared memory, and sufficient control over instruction ordering.

With \AttackName, we develop this concept into a practical, generic, and near-universally applicable probing primitive for microarchitectural latencies.
We distinguish between a \emph{padded} and an \emph{unpadded} variant.
For sufficiently large latency differences, the unpadded variant works without prior calibration, which is typically required for timer-based primitives.
For smaller latency differences, padded \AttackName adds a calibrated delay to \tprobe.
This moves the race boundary in the same way that a threshold moves the decision boundary in a timer-based attack, but without ever reading a clock.

We evaluate \AttackName on \PlatformCount CPUs spanning x86\_64, AArch64, RISC-V, and LoongArch, and \num{45} platform configurations once hybrid architectures are separated.
This breadth allows us to demonstrate \AttackName's platform-agnostic properties:
For instance, unpadded \AttackName reliably separates L1 hits from DRAM accesses on most platforms, and padded \AttackName achieves this on all of them.
This is despite vastly different core microarchitectures, interconnect topologies, and even instruction sets.
A large set of test platforms also enables us to characterize \AttackName's performance broadly.
We identify the quantity that bounds the resolution of unpadded \AttackName -- the synchronization latency \rdelta{sync} between the two threads -- and measure it across our evaluation set.
This allows us to estimate the minimum latency difference for which unpadded \AttackName works.
By quantifying the jitter in \rdelta{sync} between runs, we can additionally estimate this for the padded variant.

The results show that \AttackName enables timerless variants, e.g., of \FlushReload~\cite{Yarom2014Flush} and \EvictReload~\cite{Lipp2016} on most platforms, a conclusion we also verify experimentally.
Padded \AttackName in particular features a high temporal accuracy, and can even distinguish adjacent cache levels on some target platforms.
When multiple measurements of a target latency are permitted, we show that \AttackName can reliably resolve timing differences at the level of individual clock cycles.
We demonstrate this by recovering an address's last-level-cache (LLC) slice on a \num{12}-slice Intel CPU with \SI{96.9}{\percent} accuracy.

\AttackName is not limited to native code.
We implement \AttackName in WebAssembly and evaluate it in Chromium for x86\_64, using only the standard \code{WebWorker} parallelism and \code{SharedArrayBuffer} memory.
We find that the synchronization latency strongly correlates with the native implementation (Pearson $r = \num{0.930}$ across \num{50} matched configurations), and directly verify that browser-based \AttackName can distinguish L1 hits from DRAM accesses on all tested machines.
This indicates that \AttackName's performance in the browser approaches that of the native implementation, making it a strong primitive for browser-based attacks.
Additionally, we show that \AttackName remains effective when implemented in Lua, achieving near-perfect classification of cached and uncached accesses despite JIT compilation.


We further compare \AttackName with counter-thread timers, a common fallback when explicit timers are unavailable.
Like counter threads, \AttackName remains affected by thread placement and cache-coherence behavior.
Thus, its advantage is not independence from these mechanisms, but that it directly encodes the microarchitectural state into an architectural value instead of constructing a clock from periodically visible counter updates, which requires a (potentially adaptive) threshold.
Across physical cores, this direct race usually provides a finer decision boundary and matches or slightly exceeds the classification accuracy of counting threads.
Padded \AttackName provides the broadest coverage, while SMT-sibling counting threads remain competitive for smaller LLC-to-DRAM latency differences.

Finally, we show that \AttackName substitutes for a timer in vastly different scenarios, including complete attacks.
We adapt a first-round AES T-table attack by replacing the timed reload step with \AttackName.
On an Apple M1, where unprivileged code sees only a \SI{24}{\mega\hertz} counter, we construct minimal L2 eviction sets and run an \EvictReload-based covert channel entirely without a high-resolution timer.
These case studies demonstrate that \AttackName is a practical primitive in environments where precise timers are restricted.
Furthermore, we show the primitive is not specific to caches by detecting TLB residency, AVX wakeup throttling, and subnormal floating-point penalties.

\textbf{Contributions.} Our contributions can be summarized as follows:
\begin{compactitem}
  \item We present \AttackName, a generic timerless primitive that converts microarchitectural latency differences into the architectural outcome of a shared-memory race, and implement it on x86\_64, AArch64, RISC-V, and LoongArch.

  \item We evaluate synchronization precision, latency classification, and counter-thread baselines on \PlatformCount systems across four ISAs, showing that \AttackName outperforms counting threads in many scenarios.

  \item We implement \AttackName in WebAssembly and demonstrate its efficacy inside the browser sandbox.

  \item We demonstrate that \AttackName can replace timed probing in practical settings, including an AES T-table attack and a timerless \EvictReload-style case study.
\end{compactitem}




\section{Background}\label{sec:background}
In this section, we provide the background required for the remainder of the paper.
We first introduce software-based cache side channels and the role of memory-access latency (\Cref{sec:background:cache-attacks}).
We then discuss timers and existing timer alternatives (\Cref{sec:background:timers}).
Finally, we briefly discuss shared-memory and cache-coherence behavior (\Cref{sec:background:coherence}).

\subsection{Cache Attacks}\label{sec:background:cache-attacks}
Modern CPUs use caches to compensate for high memory access latencies.
When a program accesses a memory address, the access is fast if the corresponding cache line is in a nearby cache, and slower if the line must be fetched from a slower cache level or main memory.
Although software does not see this directly, it affects execution time.
Cache attacks exploit this timing difference to infer information about another computation~\cite{Osvik2006,Yarom2014Flush,Gruss2016Flush,Rauscher2025Systematic}.
A typical cache attack consists of three phases.
First, the attacker prepares the cache state.
Second, the victim executes and may access memory depending on secret data.
Third, the attacker probes the cache state and infers whether the victim accessed a particular cache line.
The probing phase is usually implemented as a timing measurement: the attacker accesses an address and classifies the access as a cache hit or miss based on its latency. %
Thus, cache attacks observe secret-dependent cache state and translate it into information about the victim.

\paragrabf{\FlushReload}
\FlushReload relies on shared memory between attacker and victim~\cite{Yarom2014Flush}.
The attacker flushes a shared cache line from the cache hierarchy, waits for the victim, and reloads the line.
A fast reload indicates that the victim accessed the line in the meantime.
\FlushReload is precise because it monitors individual cache lines, but it requires either a flush instruction or another mechanism to evict the target line.

\paragrabf{\PrimeProbe and \EvictReload}
\PrimeProbe does not require shared memory.
Instead, the attacker fills a cache set with its own addresses, lets the victim execute, and then measures whether its own addresses were evicted~\cite{Osvik2006}.
\EvictReload combines the reload-based observation of \FlushReload with eviction sets instead of explicit flush instructions~\cite{Lipp2016}.
This makes it applicable on systems where flushing is unavailable or restricted, but it requires finding eviction sets that map to the same cache resources as the target address.


\subsection{Timers and Timer Alternatives}\label{sec:background:timers}
In the context of side-channel attacks, a timer is typically an incrementing counter that lets software compare the duration of two operations.
Native attackers can often use architectural timers, performance counters~\cite{Lipp2016}, or cycle counters~\cite{Gerlach2023A}, such as \code{rdtsc} on x86.
In constrained environments, such as browsers, precise timers are commonly restricted or coarsened because they enable microarchitectural side channels~\cite{Schwarz2017Timers,VanGoethem2015}.
Attackers use amplification and repeated measurement to address two distinct limitations of imprecise timers.
Amplification increases the latency difference until it exceeds the timer's quantization interval~\cite{Allan2016,Schwarzl2021dynamic,Purnal2023showtime,Katzman2023gates}.
Statistical aggregation instead reduces random measurement noise when the target permits repeated observations~\cite{Bernstein2005,Schwarz2019netspectre}.
Averaging cannot, by itself, recover a latency difference that deterministically maps to the same timer value.
It is only useful when noise or sampling jitter provides information across observations.
Both approaches increase the number of required operations or victim interactions, which is problematic when the interaction is short, rate-limited, or difficult to repeat.

\paragrabf{Counting Threads}
Counting threads~\cite{Schwarz2017Timers,Lipp2016,Gras2017aslr,Hetterich2022Branch} are a common fallback in microarchitectural attacks when hardware or the execution environment lacks a precise timer.
One thread continuously increments a shared variable, while the measuring thread reads this variable before and after the operation of interest.
The observed counter difference acts as a software clock.
Its effective resolution depends not only on the counter loop, but also on how quickly updates become visible to the measuring thread.
Consequently, thread placement, cache coherence, scheduler activity, contention, and CPU frequency changes can alter the relationship between counter increments and elapsed time.

\paragrabf{Event-Based Primitives}
Some prior attacks avoid conventional timers by observing other events.
For instance, weak coherence between data and instruction caches can reveal cache activity on some platforms~\cite{Guanciale2016Cache,Jin2025Ghostcache,Thomas2025Exfilstate,thomas2026riscy}.
On some architectures, this is also possible with specialized instructions such as the Intel TSX extensions~\cite{Disselkoen2017prime,kim2022dprime,Hornetz2026tdxray} or \texttt{umwait}~\cite{Zhang2023MWAIT}.
These approaches show that explicit timestamps are not strictly necessary for side-channel attacks, but typically rely on specific ISA features, processor extensions, or execution environments.

\subsection{Cache Coherence}\label{sec:background:coherence}

Multicore CPUs allow multiple cores to access the same memory location concurrently.
To keep memory coherent, CPUs implement a cache-coherence protocol that coordinates cache-line ownership~\cite{papamarcos1984low,hum2005forward,AMD64_vol2}.
A core that wants to write to a cache line typically has to obtain exclusive ownership of that line.
If the line is not already present in a writable state, the store can trigger coherence traffic, referred to as a read-for-ownership request~\cite{Intel_combined}.
Thus, even a cached store can have a variable latency depending on the initial cache state of the written cache line.
For a single memory location, coherence establishes an order in which writes become visible.
If two threads write different values to the same location, the final value observed in memory corresponds to the write that becomes visible last in this coherence order.
This final value is architecturally observable, although the exact write ordering may depend on microarchitectural effects such as cache state, interconnect latency, and core placement.
Modern CPUs may also reorder memory operations internally.
Out-of-order execution, store buffers, and weak memory models can cause loads and stores to be issued out of program order unless the implementation uses appropriate dependencies or fences~\cite{Sorin2011}.

\section{The \AttackName Primitive}\label{sec:primitive}



In this section, we introduce \AttackName in two variants.
We first present the unpadded primitive and derive its correctness condition (\Cref{subsec:primitive-overview}).
We then introduce a padded variant for smaller latency differences (\Cref{subsec:padded-probing}), discuss the sources of variability in the race (\Cref{subsec:sources-of-imprecision}), and describe our implementation and attacker requirements (\Cref{subsec:implementation,subsec:threat-model}).

\subsection{High-Level Design}\label{subsec:primitive-overview}
\AttackName uses two concurrent threads and a shared result word $w$.
The \emph{probing thread} \tprobe executes the target operation, whose latency is either $\Delta_{fast}$ or $\Delta_{slow}$.
The \emph{reference thread} \tref executes a fixed reference operation with latency $\Delta_{ref}$.
After their respective operations complete, both threads write distinct values to $w$.
Cache coherence serializes these writes, and the value that remains in $w$ identifies the store that became visible last.
By making each store data-dependent on the preceding operation, this architectural value reveals which execution path completed later.

\begin{figure*}[t]
  \centering
  \input{figures/overview_constants.tikz}
  \begin{subfigure}{0.45\textwidth}
    \centering
    \input{figures/overview_fast.tikz}
    \subcaption{Fast target: \tref writes last, yielding $w=1$.}
  \end{subfigure}
  \hfill
  \begin{subfigure}{0.45\textwidth}
    \centering
    \input{figures/overview_slow.tikz}
    \subcaption{Slow target: \tprobe writes last, yielding $w=0$.}
  \end{subfigure}
  \caption{Unpadded \AttackName. \tprobe releases \tref through the shared synchronization word $s$ and immediately executes the target operation. \tref observes the release after $\Delta_{sync}$ and executes the reference operation. In our memory-access instantiation, the target and reference addresses contain distinct marker values, $0$ and $1$, which are propagated to the final stores. The value remaining in $w$ therefore identifies the path that completed later.}
  \label{fig:primitive}
\end{figure*}

\paragrabf{Race Construction}
\Cref{fig:primitive} illustrates the race.
Initially, \tref spins on a shared synchronization word $s$.
\tprobe starts a measurement by writing to $s$ and immediately executing the target operation with latency $\Delta_x$.
\tref observes the write after the synchronization latency $\Delta_{sync}$ and then executes its reference operation with $\Delta_{ref}$.
Ignoring jitter and assuming marker values `0' and `1', the architectural result is
$$
  w =
  \begin{cases}
    1, &\Delta_x < \Delta_{sync} + \Delta_{ref},\\
    0, &\Delta_x > \Delta_{sync} + \Delta_{ref}.
  \end{cases}
  \label{eq:photofinish-outcome}
$$
In practice, the outcome of measurements close to equality can be influenced by system-noise-induced jitter in $\Delta_{sync}$, as evaluated in \Cref{subsec:delta-sync}.

\paragrabf{Correctness Condition}
The race distinguishes the fast and slow states whenever the reference path finishes between them, i.e., $\Delta_{fast} < \Delta_{sync} + \Delta_{ref} < \Delta_{slow}$.
Unless otherwise stated, we always choose the reference operation to match the known fast case, i.e., $\Delta_{ref}=\Delta_{fast}$, as illustrated in \Cref{fig:primitive}.
The left inequality holds by construction, and the remaining requirement is $\Delta_{slow} - \Delta_{fast} > \Delta_{sync}$.
Hence, $\Delta_{sync}$ determines the temporal resolution of unpadded \AttackName: smaller synchronization latency permits the primitive to separate smaller latency differences.
Inducing a reference delay of $\Delta_{fast}$ is often trivial in practice.
With memory-access probing, for instance, both the target's fast state and the reference address can be prepared as L1 hits.
Unlike a timer-based classifier, unpadded \AttackName does not require a platform-specific numerical threshold.
The reference operation provides the comparison point directly, and simultaneous execution makes the result a relative ordering rather than an elapsed-time measurement.
We show in \Cref{sec:evaluation} that the latency difference between an L1 hit and a DRAM access exceeds $\Delta_{sync}$ on almost all tested platforms.
While we focus primarily on memory access timings, \AttackName also applies to other microarchitectural latencies (cf. \Cref{subsec:other-latency-types}), and even architectural latencies like a function's execution time.

\takeaway{Unpadded \AttackName classifies a latency without a timer or numerical threshold whenever $\Delta_{slow}-\Delta_{fast}>\Delta_{sync}$.}



\subsection{Padded \AttackName}\label{subsec:padded-probing}

\begin{figure}[t]
  \input{figures/overview_pad_constants.tikz}
  \centering
  \input{figures/overview_pad.tikz}
  \caption{Unpadded and padded \AttackName. If the target latency difference is smaller than $\Delta_{sync}$, both target states finish before the reference path. Adding a controlled delay $\Delta_{pad}$ to \tprobe shifts the race boundary so that only the fast target finishes first.}
  \label{fig:primitive_padded}
\end{figure}

Unpadded \AttackName cannot separate two target states when both finish before the reference path.
This occurs, for example, when distinguishing adjacent cache levels whose latency difference is smaller than $\Delta_{sync}$.
Padded \AttackName addresses this limitation by adding a controlled delay $\Delta_{pad}$ to \tprobe before its final store, as illustrated in \Cref{fig:primitive_padded}.
A padding value is valid if $\Delta_{pad}+\Delta_{fast} < \Delta_{sync}+\Delta_{ref} < \Delta_{pad}+\Delta_{slow}$.
Assuming no jitter, such a value exists for every non-zero latency difference.
Thus, padding moves the decision boundary like a numerical threshold in a conventional timer, but without constructing or reading a clock.
We calibrate $\Delta_{pad}$ using known instances of the fast and slow states.
Starting from zero, we increase the padding until \tprobe wins the fast case while \tref wins the slow case.
A binary search can locate this interval efficiently: if \tref wins both cases, the padding is too large; if \tprobe wins both, it is too small.
The calibration depends on the platform and the target latency pair, but not on an external timing source.
We place the padding in \tprobe rather than replacing the reference operation with a calibrated delay.
A reference path cannot complete earlier than $\Delta_{sync}$ after the release.
Thus, a delay-only reference cannot distinguish two target latencies that are both shorter than $\Delta_{sync}$.
Padding \tprobe shifts the effective boundary below $\Delta_{sync}$ to cover this case as well.
In practice, the minimum distinguishable difference is bounded by jitter in $\Delta_{sync}$, $\Delta_{fast}$, and $\Delta_{slow}$.
However, only synchronization jitter is specific to \AttackName.
Variability in target latencies also affects traditional threshold-based timing classifiers.
We quantify the synchronization jitter and the resulting uncertain range in \Cref{subsec:delta-sync}.

\takeaway{Padded \AttackName relies on one-time calibration to separate latency differences smaller than $\Delta_{sync}$.}

\subsection{Controlling Jitter}\label{subsec:sources-of-imprecision}

There are three sources of timing jitter in \AttackName: the target and reference operations, the synchronization from \tprobe to \tref, and the coherence arbitration of the final writes to $w$.
The first is inherent to the probed event and also affects timer-based classifiers.
The latter two are specific to the race and can be reduced through careful state preparation and instruction ordering.

\paragrabf{Memory and Coherence State}
We place the synchronization word $s$ and result word $w$ on separate cache lines.
With that, releasing \tref through $s$ does not share a coherence transaction with the final write race on $w$.
Before each measurement, both threads touch the two lines, placing them in their private caches and normalizing the access path.
We additionally establish the same clean Shared (\corS) state on both cores, ensuring that either final writer must obtain ownership before transitioning the result line to Modified (\corM).
This preparation requires additional care on protocols with owned dirty sharing, such as MOESI~\cite{AMD64_vol2}.
If one core holds a line in Owned (\corO) while the other holds it in Shared (\corS), the owner may upgrade to Modified with less coherence traffic than the sharer, biasing the race.
Thus, we flush or evict $s$ and $w$ to the point of coherence before both threads reload them, ensuring that neither core retains ownership.
This step is unnecessary for MESI~\cite{papamarcos1984low} and MESIF~\cite{hum2005forward}, where loading the line on both cores already establishes a clean, non-owned shared state.

\paragrabf{Instruction Ordering}
We implement the latency-critical paths in assembly and connect the relevant operations through true data dependencies.
Otherwise, compiler transformations and out-of-order execution can move work across the race boundaries.
For memory-access probing, the value loaded from the target or reference address is itself propagated to the final store to $w$.
For more general operations, we carry a marker value through a dependency chain.
Note that integrating store operations into such a dependency chain is typically not possible, as they usually do not have an output register.
Hence, we implement stores with an exchange instruction if available on the target instruction set.

\paragrabf{Thread Placement and System Activity}
$\Delta_{sync}$ depends on the placement of \tprobe and \tref, including whether they share a physical core, cache hierarchy, interconnect path, or core complex.
Hence, scheduler migration would change the decision boundary during an experiment.
Our native implementation pins both threads to fixed logical cores, and we evaluate the effect of different placements explicitly in \Cref{subsec:delta-sync}.


\subsection{Implementation}\label{subsec:implementation}

\begin{figure}[t]
  \begin{minipage}[t]{0.48\linewidth}
    \centering
    \small\textbf{\tprobe: Probing Thread}
    \begin{lstlisting}[language={},numbers=none,frame=single,
      basicstyle=\scriptsize\ttfamily,
      commentstyle=\scriptsize\itshape\color{gray!60!black},
      morekeywords={await,while},
      keywordstyle=\bfseries\color{blue!60!black},
      columns=fullflexible,breaklines=false,
      xleftmargin=12pt,xrightmargin=4pt,
      framexleftmargin=0pt,framexrightmargin=0pt,
      numbers=left,numbersep=5pt,
      numberstyle=\tiny\color{gray}]
prepare_shared(s, w)

# pre-synchronize
pre_s = 1
await pre_s == 0

# execute target
s = 1

r = target_op()   # marker 0
w = r
    \end{lstlisting}
  \end{minipage}
  \hspace{1em}
  \begin{minipage}[t]{0.43\linewidth}
    \centering
    \small\textbf{\tref: Reference Thread}
    \begin{lstlisting}[language={},numbers=none,frame=single,
      basicstyle=\scriptsize\ttfamily,
      commentstyle=\scriptsize\itshape\color{gray!60!black},
      morekeywords={await,while},
      keywordstyle=\bfseries\color{blue!60!black},
      columns=fullflexible,breaklines=false,
      xleftmargin=0pt,xrightmargin=0pt,
      framexleftmargin=0pt,framexrightmargin=0pt,
      numbers=left,numbersep=5pt,
      numberstyle=\tiny\color{gray}]
prepare_shared(s, w)
while (pre_s == 0) {
  prime(reference)
  access(s)
  access(w)
}
# synchronize
pre_s = 0
await s == 1
r = reference_op() # marker 1
w = r
    \end{lstlisting}
  \end{minipage}
  \caption{High-level pseudocode for unpadded \AttackName. \code{prepare\_shared} normalizes the cache and coherence state of $s$ and $w$; \code{target\_op} and \code{reference\_op} propagate distinct marker values through data-dependent result stores. We omit the reset logic between measurements for clarity.}
  \label{fig:pseudocode-unpadded}
\end{figure}

\Cref{fig:pseudocode-unpadded} shows the implementation at a high level.
The reference thread continuously prepares the reference operand and the shared state while waiting for a request from \tprobe.
A short pre-synchronization handshake then moves \tref into a tight spin loop on $s$, reducing variation in its response path.
After \tprobe writes to $s$, both threads execute their respective operations and propagate distinct marker values to $w$.
A separate completion flag (omitted here) ensures that \tprobe reads $w$ only after \tref has issued its result store.
For native experiments, we implement both threads entirely in assembly and expose bindings to higher-level languages.
We provide implementations for x86\_64, AArch64, RISC-V64, and LoongArch64.
The implementations differ primarily in assembly syntax and memory-ordering details.
The primitive itself requires only loads and stores, concurrent execution, and shared memory.
For padded \AttackName, we generate $\Delta_{pad}$ with a chain of dependent register-only additions whose final value feeds the result store.
Varying the chain length provides fine-grained control over the padding without relying on a timer.




\subsection{Threat Model}\label{subsec:threat-model}

We assume an attacker with unprivileged code execution on the target machine.
\AttackName inherently requires (i) two threads that can execute concurrently, (ii) writable shared memory for synchronization, completion, and the result word, and (iii) the ability to invoke the target and reference operations.
It does not require a timestamp source, performance counter, privileged instruction, or ISA-specific extension.
Our evaluated native implementation also assumes the attacker can keep the two threads on stable logical cores, e.g., through the affinity interface available to unprivileged users on Linux and Windows.
Without explicit affinity control, scheduler migration increases $\Delta_{sync}$ variability and may require repeated measurements, but this does not change the construction.
Our padded implementation generates dependency chains at runtime for convenience.
However, runtime code generation is not necessary: pre-generated chains or a dependency-preserving loop can provide the same controlled delay.



\section{Evaluation}\label{sec:evaluation}

We evaluate four properties of \AttackName: the synchronization precision that bounds its temporal resolution (\Cref{subsec:delta-sync}), its single-shot accuracy when classifying cache states (\Cref{subsec:measuring-memory-access-latencies}), its performance relative to counting-thread timers (\Cref{subsec:comparison-to-counter-threads}), and its applicability inside a web browser (\Cref{subsec:attackname-in-the-browser}).
Our test set contains \PlatformCount physical CPUs spanning x86\_64, AArch64, RISC-V64, and LoongArch64.
Seven of these CPUs contain heterogeneous core types.
Treating each core type as a separate platform yields \num{45} native platform configurations.
\Cref{tab:precision,tab:precision_ht,tab:precision_wasm} in \Cref{sec:precision-analysis} report full per-platform results.
Unless stated otherwise, we pin \tprobe to logical core 0 and \tref to the next physical core.
On the \num{20} configurations supporting simultaneous multithreading (SMT), we repeat each experiment with the threads on sibling logical cores.
Where possible, we disable hardware prefetchers.
We express short delays in \emph{\texttt{add}-cycles}, defined as the latency of one register-only \texttt{add} in a dependency chain.
This unit is portable across our four ISAs and avoids conflating core cycles with fixed-frequency timestamp-counter ticks.
On many microarchitectures, one \texttt{add}-cycle corresponds to one core cycle~\cite{abel2019uops,Intel_combined}.


\subsection{Synchronization Precision}\label{subsec:delta-sync}

The resolution of unpadded \AttackName depends on the synchronization latency \rdelta{sync} and, more importantly for padded \AttackName, on its run-to-run jitter.
Thus, we measure both with respect to thread placement.

\paragrabf{Method}
To measure \rdelta{sync}, we instantiate a padded race without target or reference memory accesses.
After \tprobe releases \tref through the synchronization word, \tprobe first executes a dependency chain of $d$ register-only \texttt{add} instructions, while \tref writes after observing the release.
Thus, the result compares the controlled padding $d$ against \rdelta{sync}.
We execute \num{2048} such races for each $d \in [0, 1000)$ and record the probability $p(d)$ that \tprobe wins.
At $d=0$, \tprobe normally wins, and as $d$ increases, the winner transitions to \tref.
On most CPUs, this transition zone consists of multiple $d$ values where both threads have a high chance of winning.
This is illustrated in \Cref{fig:sync-time}.
We call this the \emph{uncertain range} $U$, and define the lower endpoint as the largest $d$ for which $p(d)\geq\SI{95}{\percent}$, and its upper endpoint as the smallest $d$ for which $p(d)\leq\SI{5}{\percent}$, i.e., $U$ brackets the transition.
Its midpoint estimates \rdelta{sync}, while its width captures the observable synchronization jitter.
The upper edge of $U$ provides a conservative bound for unpadded \AttackName, whose condition is $\Delta_{slow}-\Delta_{fast}>\Delta_{sync}$.
For padded \AttackName, the width $|U|$ approximates the minimum latency separation that can be classified reliably with one race.
Smaller differences can still be recovered by combining repeated measurements, as we demonstrate in \Cref{subsec:determining-llc-slices}.
We run this experiment on all \num{45} platform configurations.
To study topology effects, we keep \tprobe on core 0 and place \tref on every other logical core.
All reported values of \rdelta{sync} are midpoints of the corresponding uncertain ranges.

\begin{figure}[t]
  \flushleft
  \input{plots/sync_time_common.tikz}
  \input{plots/sync_time_Intel-Core-i7-8700.tikz}
  \input{plots/sync_time_AMD-EPYC-7252.tikz}
  \input{plots/sync_time_Loongson-3a600.tikz}
  \caption{Synchronization transition on representative CPUs ($n=2048$ races per padding value). Increasing the dependent padding in \tprobe moves the winner from \tprobe to \tref. The transition midpoint estimates \rdelta{sync}. The shaded uncertain range captures values for which neither thread wins at least \SI{95}{\percent} of races. Core placement affects both quantities.}
  \label{fig:sync-time}
\end{figure}

\paragrabf{Cross-Core versus SMT Placement}
SMT siblings provide a shorter synchronization path because both threads share the private cache hierarchy.
However, they also share frontend and execution resources, and we observe a substantially wider uncertain range on most SMT configurations.
Thus, SMT reduces the mean synchronization latency but generally worsens single-shot temporal precision.
For the remainder of the paper, we prefer separate physical cores unless stated otherwise.
Across separate physical cores, the exact core pair has little effect on many CPUs.
The measured \rdelta{sync} is identical across placements on \num{8} platforms and varies by fewer than \num{10} \texttt{add}-cycles on another \num{21}.
The main exception is topology that crosses an internal interconnect boundary.
For example, on the AMD \labthirtyCPU, \rdelta{sync} is \num{78} \texttt{add}-cycles when both threads reside in the same core complex (CCX), but averages \num{563} \texttt{add}-cycles across CCXs.
We observe the same behavior on all tested multi-CCX AMD systems: selecting cores within one CCX significantly improves the race boundary.
On several platforms, $U$ is only one \texttt{add}-cycle wide, the finest granularity our setup resolves.
Repeated experiments on these systems distinguish chains differing by one dependent \texttt{add}.
While this result characterizes the synchronization path, variability in the target operation can widen the effective transition during an attack.
Accordingly, the values in \Cref{sec:precision-analysis} are optimistic lower bounds on the total classification noise.

\paragrabf{Predicting L1-versus-DRAM Feasibility}
Knowledge of \rdelta{sync} enables predicting the success of, e.g., a cache-state classifier.
To demonstrate this, we sweep \rdelta{pad} from \num{0} to \num{1000} \texttt{add}-cycles for both an L1 hit and a DRAM access, taking \num{2048} samples per case and padding value.
We conservatively estimate the difference between an L1 hit and a DRAM access $\Delta^{L1}_{DRAM}$ as the separation between the two transition regions: the highest padding for which the L1 case still wins more than \SI{5}{\percent} of races and the lowest padding for which the DRAM case wins fewer than \SI{95}{\percent}.
When running on core 0 and a separate physical core, $\Delta^{L1}_{DRAM}$ exceeds \rdelta{sync} on all but two platforms, \labfourtyfourCPU and \labeightythreeCPU.
The unpadded condition therefore predicts successful L1-versus-DRAM classification on the remaining platforms.
In contrast, $\Delta^{L1}_{DRAM}$ exceeds $|U|$ on every platform, predicting that calibrated padding can separate the two states throughout the test set.
The direct classification experiments below, together with the AES case study in \Cref{subsec:aes-t-table}, confirm the padded prediction on every platform.
Both platforms predicted to fail without padding indeed do, although two further platforms fall short for reasons we discuss below.

\subsection{Cache-Latency Classification}\label{subsec:measuring-memory-access-latencies}

In this section, we evaluate \AttackName as a primitive for inferring a cache line's state.
We consider L1 hits versus DRAM accesses, common in \FlushReload-style attacks~\cite{Yarom2014Flush}, and LLC hits versus DRAM accesses, as required by LLC \PrimeProbe~\cite{Liu2015Last}.
For padded \AttackName, we additionally test the smaller L1-versus-L3 and L1-versus-L2 differences.

\paragrabf{Method}
For every combination of cache state for the address and padding value, we collect $2^{14}$ samples from each class, with \rdelta{pad} from \num{0} to \num{500} \texttt{add}-cycles.
Note that padding zero is the unpadded construction.
We report the $F_1$ score with the fast state as the positive class.
We prepare cache levels using physical-address-based eviction sets, as detailed in \Cref{sec:reliable-cache-eviction}.
Where a sufficiently precise timer is available, we independently verify that eviction leaves the target at the intended cache level.
We cannot perform this validation on the RISC-V and LoongArch64 systems or on \labseventyoneCPU.
Thus, their L2/L3 results are conservative: a low score may indicate imperfect state preparation rather than insufficient resolution of \AttackName.
For LLC-versus-DRAM measurements on SMT siblings, \tref also accesses an LLC-resident reference line.
With separate physical cores, we instead retain an L1-resident reference because \rdelta{sync} already shifts the reference path beyond typical LLC latency.
Avoiding an LLC access in \tref removes interconnect-induced jitter~\cite{paccagnella2021lord} and improves the decision boundary.

\paragrabf{Unpadded \AttackName}
For L1 versus DRAM on separate physical cores, \num{41} of \num{45} platforms achieve $F_1\geq0.97$.
The four exceptions expose two distinct failure modes.
On \labseventyoneCPU and the efficiency cores of \labninetyfiveCPU, DRAM latency varies enough that some nominally slow probes complete before the reference path.
On \labfourtyfourCPU and \labeightythreeCPU, the L1-to-DRAM separation is too small relative to \rdelta{sync}, violating the unpadded condition.
Among all configurations, \num{38} produce fewer than \num{20} errors across \num{32768} classified accesses, and \num{4} classify every access correctly.
All tested RISC-V and LoongArch64 platforms achieve high scores, showing that the unpadded primitive supports timer-free, calibration-free \FlushReload-style attacks beyond x86 and ARM.
On SMT siblings, all \num{20} platforms exceed $F_1=0.99$ for L1 versus DRAM, with \num{18} platforms having fewer than \num{20} errors and \num{8} platforms being perfect.
The shorter synchronization path dominates the additional SMT jitter for this large latency difference.
For LLC versus DRAM, the cross-core configuration again performs strongly: \num{37} of \num{45} platforms exceed $F_1=0.99$, including \num{5} perfect classifiers.
These platforms support an unpadded, calibration-free LLC \PrimeProbe readout.
On SMT siblings, \num{6} of \num{20} platforms exceed $F_1=0.99$.
Here, the smaller cache-latency separation is more sensitive to the wider uncertain range, confirming the placement result from \Cref{subsec:delta-sync}.

\paragrabf{Padded \AttackName}
Padding converts a single race boundary into a tunable classifier.
\Cref{fig:padding-results} illustrates how padding allows tuning the classifier on a sample \labthirtythreeCPU.
At zero padding, both the fast and slow target may finish before \tref, resulting in an unusable classifier with $F_1\approx 2/3$.
As padding increases, the slow state crosses the reference boundary first and the fast state later.
Between these transitions, both classes are classified accurately.
We define the \emph{golden range} $G$ as a contiguous interval of padding values for which $F_1>0.98$.
Its width measures calibration tolerance: a wider $G$ makes the classifier less sensitive to small shifts in execution conditions.
Every platform has a non-empty golden range for L1 versus DRAM\@.
Thus, padded \AttackName provides timer-free \FlushReload-style classification on all tested CPUs, including platforms where the unpadded condition fails.
For LLC versus DRAM, \num{42} of \num{45} platforms have a golden range.
The finer L1-versus-L3 difference is separable with one race on \num{32} platforms, and even L1 versus L2 is separable on \num{15}.
The latter includes the efficiency cores of \labseventytwoCPU and the Apple M1.
On the M1, the achieved granularity significantly exceeds that of the unprivileged architectural timer~\cite{ARM2021GenericTimer,Lipp2016}.
With SMT placement, all \num{20} platforms retain a golden range for L1 versus DRAM.
Smaller differences are less robust: \num{13} platforms expose a golden range for L1 versus L3, and \num{7} for LLC versus DRAM.
This confirms the tradeoff established in \Cref{subsec:delta-sync}, i.e., SMT minimizes \rdelta{sync}, whereas separate physical cores usually minimize the transition width.

\begin{figure*}[t]
  \centering
  \input{plots/optimizer_lab33_grouped.tikz}
  \caption{Cache-state classification on \labthirtythreeCPU with cross-core placement ($n=16384$ per class and padding value). The fast-state accuracy $A_{fast}$ decreases at a larger padding value than the slow-state accuracy $A_{slow}$. Their overlap forms the golden range $G$, in which the combined classifier achieves $F_1>0.98$.}
  \label{fig:padding-results}
\end{figure*}

\takeaway{Unpadded \AttackName provides calibration-free L1-versus-DRAM classification on \num{41} of \num{45} cross-core platforms. Calibrated padding extends coverage to all tested CPUs and can even classify adjacent cache levels.}


\subsection{Comparison with Counting Threads}\label{subsec:comparison-to-counter-threads}

Counting threads are the standard timing fallback when an explicit high-resolution timer is unavailable~\cite{Wray1992,Lipp2016,Schwarz2017Timers,Gras2017aslr,Schwarz2020malware,Hetterich2022Branch}.
They require the same basic resources as \AttackName, i.e., parallel execution and shared memory, but expose time through periodically visible counter updates.
A cache attack must then calibrate a numerical threshold over this emulated clock.
We compare the two approaches first at the level of temporal resolution and then through direct cache-state classification.

\paragrabf{Counting-Thread Implementation}
Our assembly implementation is a loop that increments a register and stores it to shared memory, requiring three instructions per iteration.
We evaluate all \num{45} native platform configurations cross-core and the \num{20} SMT configurations on sibling threads.

\paragrabf{Effective Counter Resolution}
The relevant resolution is not the local loop throughput but the interval between counter values observable by another core.
Store buffering, cache ownership, and coherence propagation can combine or delay updates.
We call the mean interval between globally observed changes the \emph{edge width}.
A branch-free observer repeatedly loads the counter and counts value changes.
We measure the wall-clock time required to observe \num{10000000} changes using a coarse OS timer, then convert it to \texttt{add}-cycles by benchmarking \num{65536} repetitions of a \num{1024}-instruction dependent-add chain.
The observer introduces a small positive bias, but the counter loop itself takes fewer than \num{7} \texttt{add}-cycles per iteration on every platform, whereas the measured edge width is substantially larger on nearly all of them.
The cross-core edge width is higher than \rdelta{sync} by more than \num{30} \texttt{add}-cycles on \num{28} of \num{45} platforms, while the opposite occurs on only \num{5}.
Thus, the unpadded race usually has a finer boundary than a cross-core counting thread.
On SMT siblings, edge width and \rdelta{sync} are typically within \num{15} \texttt{add}-cycles, although \labninetyoneCPU and \labseventynineCPU are outliers at \num{3} and \num{249} \texttt{add}-cycles, respectively.
For padded \AttackName, the relevant quantity is $|U|$ rather than \rdelta{sync}.
On \num{19} of \num{20} platforms, the cross-core uncertain range is narrower than the edge width of an SMT-sibling counting thread, indicating a substantially finer calibrated decision boundary.

\paragrabf{Direct Classification}
We use the counting thread as a timer for the same L1-versus-DRAM and LLC-versus-DRAM classifications.
For each pair, we profile the mean fast and slow readings, set the threshold to
$T_{fast}+\frac{1}{2}(T_{slow}-T_{fast})$, and classify the fast and slow cases over \num{65536} iterations.
In the cross-core setting, counting threads exceed $F_1=0.97$ on \num{36} of \num{40} platforms for L1 versus DRAM and \num{32} for LLC versus DRAM.
On the same subset, unpadded \AttackName reaches \num{37} and \num{33}, respectively, despite requiring no numerical threshold.
Padded \AttackName reaches all \num{40} platforms for L1 versus DRAM and \num{35} for LLC versus DRAM.
Thus, cross-core \AttackName matches or slightly exceeds the counting-thread baseline, and padding further increases coverage.
The SMT result is more nuanced.
Counting threads exceed $F_1=0.97$ on \num{16} of \num{18} platforms for both latency pairs.
Unpadded \AttackName reaches all \num{18} for L1 versus DRAM but only \num{7} for LLC versus DRAM. Padded \AttackName exceeds $F_1=0.98$ on all \num{18} and \num{9}, respectively.
Consequently, an SMT-sibling counting thread remains the stronger baseline for the smaller LLC-to-DRAM difference, but at the cost of threshold calibration.


\takeaway{Across physical cores, \AttackName is at least as accurate as counting-thread timing while unpadded \AttackName eliminates threshold calibration. Padded \AttackName provides the broadest coverage.}

\subsection{Browser Execution}\label{subsec:attackname-in-the-browser}

We demonstrate that \AttackName in WebAssembly is a practical alternative to counting threads in browsers.
Browsers are an important restricted environment for microarchitectural attacks because they execute untrusted code, support parallel execution through Web Workers, and expose shared memory through \texttt{SharedArrayBuffer} under cross-origin isolation~\cite{Oren2015,Roettger2021SpectreJS,W3C2015WebWorker,FirefoxSharedArrayBuffer}.

\paragrabf{Implementation and Scope}
We implement both threads in WebAssembly and execute them in Chromium v147.0 on x86\_64.
Using JavaScript, we rewrite the WebAssembly bytecode to insert a variable-length chain of dependent additions for \rdelta{pad}.
Inspection of the generated machine code confirms that Chromium emits the intended register-only dependency chain.
The primitive itself uses only Web Workers and \texttt{SharedArrayBuffer}.
Standard browser code cannot execute \texttt{clflush}, and this experiment requires known cache states to measure classification accuracy.
Thus, we use a patched Chromium build to prepare and label the probe line through native cache-management instructions.
We do not use these hooks for synchronization, padding, or readout, and the WebAssembly primitive does not normalize the coherence state of $s$ and $w$ as the native implementation does.
We repeat the synchronization and L1-versus-DRAM experiments with the same cross-core and SMT placements as in native code, again treating heterogeneous core types separately.
Per padding value, we collect \num{2048} synchronization races and \num{4000} cache-state samples.

\paragrabf{Results}
The WebAssembly and native results are strongly correlated (Pearson $r=0.930$ across the matched configurations).
Browser execution increases the mean \rdelta{sync} from \num{98.7} to \num{129} \texttt{add}-cycles and widens the mean uncertain range from \num{23.7} to \num{37.8} \texttt{add}-cycles.
The additional latency and jitter are consistent with the longer generated response path and the missing coherence-state preparation for $s$ and $w$.
Despite this reduced resolution, unpadded WebAssembly \AttackName obtains $F_1\geq0.99$ in \num{46} of \num{50} configurations for L1 versus DRAM.
Padded \AttackName yields a golden range on all \num{50} configurations.
Its mean width is \num{84.8} \texttt{add}-cycles, close to the native mean of \num{86.9}, and the per-configuration widths remain strongly correlated (Pearson $r=0.853$).
Thus, the shared-memory race survives browser compilation with only a moderate increase in synchronization noise.

\takeaway{WebAssembly preserves the behavior and accuracy of \AttackName: padded classification succeeds in every tested browser configuration.}


\section{Case Studies}\label{sec:case-studies}
In this section, we show that \AttackName can substitute for a timer in \num{5} case studies.
\Cref{subsec:aes-t-table} demonstrates \AttackName in a timerless attack against AES\@.
In \Cref{subsec:timerless-evict-reload}, we show that \AttackName can find minimal eviction sets, and use them for an Evict+Reload-based covert channel on an Apple M1 CPU\@.
\Cref{subsec:determining-llc-slices} demonstrates that \AttackName can detect the LLC slice of an address with aggregated measurements.
\Cref{subsec:other-latency-types} shows that \AttackName also resolves non-cache latencies.
Finally, in \Cref{subsec:attackname-in-lua}, we show that \AttackName does not strictly require access to binary code by implementing the primitive in Lua.

\subsection{AES T-Table Attack}\label{subsec:aes-t-table}

\input{plots/t-table-matrix.tikz}
\begin{figure}[t]
    \centering
    \begin{tikzpicture}
        \begin{groupplot}[
            group style={group size=2 by 2, horizontal sep=1.25cm, vertical sep=1.25cm,},
        ]
        \ttablematrix{data/t-table/lab89.txt}{Intel \labeightynineCPU}
        \ttablematrix{data/t-table/rpi3b.txt}{\labninetyeightCPU}
        \ttablematrix{data/t-table/lab53.txt}{\labfiftythreeCPU}
        \ttablematrix{data/t-table/lab91.txt}{\labninetyoneCPU}
        \end{groupplot}
    \end{tikzpicture}
    \caption{A selection of heatmaps produced by the T-table attack using unpadded \AttackName on \num{4} different CPUs. All show the diagonal line expected when the first key byte is 0.}
    \label{fig:t-table-matrices}
\end{figure}

To evaluate \AttackName's performance in a real-world attack scenario, we use it in a Flush+Reload-style attack targeting OpenSSL's T-table-based AES implementation.
Specifically, we mount a first-round chosen plaintext attack to partially recover the first key byte from memory access patterns~\cite{Osvik2006}.
This attack has emerged as a standard benchmark for cache attacks, making our results directly comparable to a large body of existing work~\cite{Gruss2016Flush,Disselkoen2017prime,VanSchaik2018malicious,Briongos2020reload,Lipp2020takeaway,Purnal2021PrimeScope,Gerlach2023A,Purnal2023showtime,Gerlach2024Hash,Weber2024no,Thomas2025Exfilstate}.

\paragrabf{Setup}
We specifically target OpenSSL v.3.5.0, compiled for all \num{4} architectures in our evaluation set using GCC v15.3.
OpenSSL's T-table AES implementation uses no architecture-specific code, enabling us to mount the experiment on all \PlatformCount CPUs in our evaluation set.
However, we exclude Apple's M1 and M2 architectures, as they feature a cache line size of \SI{128}{\byte} at the point of coherence, whereas the attack requires \SI{64}{\byte}.
We run unpadded \AttackName on separate physical cores, with \tprobe on core 0 and \tref on the physical core with the next-highest ID.
As with the experiments in \Cref{sec:evaluation}, we treat CPUs with hybrid architectures as two test platforms.
With \num{5} of our \num{36} remaining CPUs being hybrids, this results in \num{41} total platforms.
Since the memory containing the T-tables is physically shared between the attacker and victim, we can use \AttackName in a Flush+Reload-like attack.
We flush each relevant cache line to DRAM, perform the encryption, and then probe the cache lines using \AttackName.
This corresponds to the L1 vs.\ DRAM configuration tested in \Cref{subsec:measuring-memory-access-latencies}.
On platforms without a flush instruction, we evict cache lines using the method described in \Cref{sec:reliable-cache-eviction}.
In each iteration of the attack, we twice encrypt a chosen 128-bit plaintext block with \num{0} as the first key byte.
We set the first byte of this block to a multiple of 16, with the rest containing random data generated from a fixed seed.
Overall, we run \num{1000} probing iterations for each possible setting of the first byte, probing all \num{16} relevant cache lines in each iteration.
We analyze the resulting access patterns to determine whether the attack is successful and repeat the attack \num{1000} times per platform to gauge its reliability.

\paragrabf{Results}
Of the \num{41} tested platforms, \num{37} have a success rate of \SI{100}{\percent}.
One platform, the \labnineCPU, reaches \SI{99.9}{\percent}.
\Cref{fig:t-table-matrices} shows a representative sample of the heatmaps the attack yields.
All clearly show the diagonal line expected when the first key byte is 0.
Moreover, the noise patterns in non-diagonal cells show similarities across machines, indicating high probing accuracy.
On average across all platforms, the median runtime is \SI{315}{\milli\second}.
The lowest average runtime is \SI{218}{\milli\second} on the AMD \labeightyCPU, and the highest is \SI{17.4}{\second} on the \labfourtyeightCPU.
There are only \num{3} machines on which the attack is unreliable: the \labthirtytwoCPU (\SI{3.5}{\percent} success rate), the \labeightythreeCPU (\SI{55}{\percent} success rate), and the Intel \labfourtyfourCPU (\SI{0}{\percent} success rate).
For the latter two, this is expected, as \Cref{subsec:measuring-memory-access-latencies} predicts: on these platforms $\Delta^{L1}_{DRAM}$ does not exceed \rdelta{sync}, so the precondition for the unpadded variant does not hold.
We confirm that a padding of \num{75} \texttt{add}-cycles restores a \SI{100}{\percent} success rate.
For the \labthirtytwoCPU, which has only two parallel threads, we believe the low success rate results from its small cache (\SI{2}{\mega\byte} LLC) combined with its high susceptibility to outside interference due to its low core count.

\subsection{Evict+Reload}\label{subsec:timerless-evict-reload}

In addition to demonstrating \AttackName in a Flush+Reload attack, we show that attackers can use it to find eviction sets without a precise timer.
For this, we implement a cross-core covert channel on an Apple M1 system using Evict+Reload on the L2 cache.
On this architecture, traditional timing-based cache attacks are challenging since only a coarse-grained \SI{24}{\mega{}Hz} timer (CNTVCT) is available.

\paragrabf{Finding Minimal Eviction Sets}
The M1 processor uses a complex hash function to index its L2 cache~\cite{Yu2023Synchronization}.
Its hash depends on physical address bits outside the page offset, so we cannot compute minimal eviction sets from virtual addresses alone.
Instead, we start with an excessive eviction set of \num{1536} addresses that have the same page offset as the victim address, \ie a \SI{24}{\mega\byte} array.
We then remove one address at a time and use the padded \AttackName primitive to detect whether the eviction set can still evict the target address from the L2 cache.
If eviction still succeeds, we remove this address from the eviction set, and repeat until we reach \num{12} addresses, i.e., the L2 cache's associativity.
While we only use minimal eviction sets in our experiment, note that minimality is not a strict requirement.
The initial set of \num{1536} could also be used with a potential reduction in throughput.

\paragrabf{Synchronization}
We use the coarse-grained \texttt{CNTVCT} timer to lay out fixed time slots for each transmitted bit.
To signal the start of a transmission, the sender holds the one-signal for eight consecutive time slots, which the receiver decodes as a start byte of \code{0xff}. The receiver discards all bytes until it observes this value.
If no coarse-grained timer is available, an additional cache line could transmit a clock signal, or post-processing could align the measurement to bits.

\paragrabf{Receiver}
The receiver uses two threads with the \AttackName primitive in a loop.
Since \rdelta{sync} is sufficiently long, we use an L1 cache hit and additional padding instructions instead of an L2 cache hit in the reference thread.
After each measurement, we evict the shared cache line from the L2 cache using a minimal eviction set.
A counter is initialized to 0 at the start of a time slot.
On each cache hit, it is incremented; on a cache miss, it is decremented.
If the counter is positive at the end of a time slot, the received value is 1; otherwise, it is 0.

\paragrabf{Sender}
To send a 1, the sender repeatedly accesses the shared cache line in a loop.
To send a 0, the sender evicts the shared cache line from its L1 cache.
This is necessary since eviction by the receiver from the L2 cache can fail if the value is still present in the L1 cache of a different core.

\paragrabf{Results}
We run the covert channel on the M1's Firestorm performance cores at \SI{30}{\kilo\byte/\second} and transmit \SI{64}{\kilo\byte} of data, conducting \num{200} runs in total.
We find a minimal eviction set in \num{163} of these runs (\SI{81.50}{\percent}) with a median construction runtime of \SI{289}{\milli\second}.
If an eviction set is found, a median of \SI{96.87}{\percent} of bits is received correctly without error correction.
This shows that the \AttackName primitive can reliably conduct eviction-based cache side-channel attacks on platforms without highly accurate timers.

\subsection{Determining LLC Slices}\label{subsec:determining-llc-slices}

\begin{figure}[t]
  \begin{tikzpicture}

\definecolor{darkcyan32144140}{RGB}{32,144,140}
\definecolor{darkgray176}{RGB}{176,176,176}
\definecolor{darkslateblue48103141}{RGB}{48,103,141}
\definecolor{darkslateblue5882139}{RGB}{58,82,139}
\definecolor{darkslateblue6857130}{RGB}{68,57,130}
\definecolor{greenyellow20222430}{RGB}{202,224,30}
\definecolor{indigo68184}{RGB}{68,1,84}
\definecolor{indigo7230112}{RGB}{72,30,112}
\definecolor{lightgray204}{RGB}{204,204,204}
\definecolor{mediumseagreen32164133}{RGB}{32,164,133}
\definecolor{mediumseagreen54184119}{RGB}{54,184,119}
\definecolor{mediumseagreen9420197}{RGB}{94,201,97}
\definecolor{teal39124142}{RGB}{39,124,142}
\definecolor{yellowgreen14621565}{RGB}{146,215,65}

\newcommand{\llcplotwidth}[0]{6.75cm}
\newcommand{\llcplotheight}[0]{3cm}

\begin{groupplot}[group style={group size=1 by 1}]
\nextgroupplot[
height=\llcplotheight, width=\llcplotwidth,
tick align=outside,
tick pos=left,
x grid style={darkgray176},
xlabel={\tprobe Padding},
xmin=0, xmax=80,
xtick style={color=black},
xlabel style={font=\small},
xticklabel style={font=\small},
y grid style={darkgray176},
ylabel={\tprobe wins (\%)},
ymin=-3, ymax=104,
ytick style={color=black},
ylabel style={font=\small},
yticklabel style={font=\small},
legend entries={0,1,2,3,4,5,6,7,8,9,10,11},
legend columns=2,
legend cell align={left},
legend style={
  at={(1.03,1)}, anchor=north west, inner ysep=-.5pt, inner xsep=2pt,
  draw=lightgray204, fill opacity=0.8, draw opacity=1, text opacity=1,
  font=\scriptsize, row sep=-4pt,
  /tikz/every even column/.append style={column sep=1pt}
},
legend image code/.code={\draw[mark repeat=2,mark phase=2,#1] plot coordinates {(0cm,0cm) (0.2cm,0cm)};},
legend image post style={line width=0.8pt}
]
\addplot [line width=0.64pt, indigo68184]
table {%
0 100
1 100
2 100
3 100
4 100
5 100
6 100
7 100
8 100
9 100
10 99.9755859375
11 100
12 100
13 100
14 100
15 100
16 100
17 99.9267578125
18 100
19 99.8046875
20 99.9755859375
21 100
22 99.951171875
23 99.951171875
24 99.951171875
25 99.8291015625
26 99.853515625
27 99.90234375
28 99.8779296875
29 99.609375
30 99.51171875
31 99.21875
32 99.462890625
33 99.7802734375
34 99.9267578125
35 99.951171875
36 99.3896484375
37 98.9990234375
38 98.7548828125
39 99.0478515625
40 98.6083984375
41 98.6083984375
42 97.94921875
43 96.3134765625
44 97.900390625
45 95.361328125
46 95.9228515625
47 93.359375
48 90.966796875
49 91.5283203125
50 90.0146484375
51 89.208984375
52 94.2138671875
53 90.087890625
54 90.3564453125
55 91.455078125
56 87.9150390625
57 90.8447265625
58 87.1826171875
59 84.47265625
60 80.0537109375
61 78.8818359375
62 79.833984375
63 77.1728515625
64 76.7822265625
65 67.724609375
66 62.5244140625
67 60.83984375
68 52.2705078125
69 38.720703125
70 40.283203125
71 30.46875
72 24.6826171875
73 22.216796875
74 19.4580078125
75 12.1337890625
76 6.6650390625
77 6.2255859375
78 5.46875
79 5.5419921875
80 4.150390625
81 2.7587890625
82 1.9287109375
83 0.732421875
84 0.9033203125
85 0.7080078125
86 2.1484375
87 0.5615234375
88 0.78125
89 0.09765625
90 0.2197265625
91 0.2685546875
92 1.220703125
93 0.09765625
94 0.2685546875
95 0.2685546875
96 1.2939453125
97 0.830078125
98 0.341796875
99 0.9521484375
100 1.318359375
101 1.4404296875
102 0.78125
103 0.8544921875
104 0.1953125
105 1.1474609375
106 1.0009765625
107 0.5126953125
108 0.9521484375
109 0.8544921875
110 0
111 0.0244140625
112 0.0732421875
113 0.48828125
114 0.78125
115 1.0009765625
116 0.3662109375
117 0.244140625
118 0.5615234375
119 0.244140625
120 0.830078125
121 0
122 0.3662109375
123 0.29296875
124 0.0244140625
125 0.1708984375
126 0.0244140625
127 0.048828125
128 0
129 0
130 0.0732421875
131 0
132 0.0244140625
133 0
134 0
135 0.0244140625
136 0.0244140625
137 0.0244140625
138 0.048828125
139 0
140 0
141 0
142 0
143 0
144 0
145 0
146 0
147 0
148 0
149 0
150 0
151 0.0244140625
152 0
153 0
154 0
155 0
156 0
157 0.0244140625
158 0
159 0
160 0
161 0
162 0
163 0
164 0
165 0
166 0
167 0
168 0.0244140625
169 0
170 0
171 0
172 0
173 0
174 0
175 0
176 0
177 0
178 0
179 0
180 0
181 0
182 0
183 0
184 0
185 0
186 0
187 0
188 0
189 0
190 0
191 0
192 0
193 0
194 0
195 0
196 0
197 0
198 0
199 0
200 0
201 0
202 0
203 0
204 0
205 0
206 0
207 0
208 0
209 0.0244140625
210 0
211 0
212 0
213 0
214 0
215 0
216 0
217 0
218 0
219 0
220 0
221 0
222 0
223 0
224 0
225 0
226 0
227 0
228 0
229 0
230 0
231 0
232 0
233 0.0244140625
234 0
235 0
236 0
237 0
238 0
239 0
240 0
241 0
242 0
243 0.0244140625
244 0
245 0
246 0
247 0
248 0
249 0
250 0
251 0
252 0
253 0
254 0
255 0
256 0
257 0
258 0
259 0
260 0
261 0
262 0
263 0
264 0
265 0
266 0
267 0
268 0
269 0
270 0
271 0
272 0
273 0
274 0
275 0
276 0
277 0
278 0
279 0
280 0
281 0
282 0
283 0
284 0
285 0
286 0
287 0
288 0
289 0
290 0
291 0
292 0
293 0
294 0
295 0
296 0
297 0
298 0
299 0
300 0
301 0
302 0
303 0
304 0
305 0
306 0
307 0
308 0
309 0
310 0
311 0
312 0
313 0
314 0
315 0
316 0
317 0
318 0
319 0
320 0
321 0
322 0
323 0
324 0
325 0
326 0
327 0
328 0
329 0
330 0
331 0
332 0
333 0
334 0
335 0
336 0
337 0
338 0
339 0
340 0
341 0
342 0
343 0
344 0
345 0
346 0
347 0
348 0
349 0
350 0
351 0.0244140625
352 0
353 0
354 0
355 0
356 0
357 0
358 0
359 0
360 0
361 0
362 0
363 0
364 0
365 0
366 0
367 0
368 0
369 0
370 0
371 0
372 0
373 0
374 0
375 0
376 0
377 0
378 0
379 0
380 0
381 0
382 0
383 0
384 0
385 0
386 0
387 0
388 0
389 0
390 0
391 0
392 0
393 0
394 0
395 0
396 0
397 0
398 0
399 0
400 0
401 0
402 0
403 0
404 0
405 0
406 0
407 0
408 0
409 0
410 0
411 0
412 0
413 0
414 0
415 0
416 0
417 0
418 0
419 0
420 0
421 0
422 0
423 0
424 0
425 0
426 0
427 0
428 0
429 0
430 0
431 0
432 0
433 0
434 0
435 0
436 0
437 0
438 0
439 0
440 0
441 0
442 0
443 0
444 0
445 0
446 0
447 0
448 0
449 0
450 0
451 0
452 0
453 0
454 0
455 0
456 0
457 0
458 0
459 0
460 0
461 0
462 0
463 0
464 0
465 0
466 0
467 0
468 0
469 0
470 0
471 0
472 0
473 0
474 0
475 0
476 0
477 0
478 0
479 0
480 0
481 0
482 0
483 0
484 0
485 0
486 0
487 0
488 0
489 0
490 0
491 0
492 0
493 0
494 0
495 0
496 0
497 0
498 0
499 0
};
\addplot [line width=0.64pt, indigo7230112]
table {%
0 100
1 99.9755859375
2 100
3 99.9755859375
4 99.9755859375
5 100
6 100
7 100
8 100
9 100
10 100
11 100
12 100
13 100
14 99.8779296875
15 100
16 99.8779296875
17 99.9267578125
18 100
19 99.9267578125
20 99.951171875
21 99.9267578125
22 99.8046875
23 99.90234375
24 99.7802734375
25 99.3896484375
26 99.853515625
27 99.755859375
28 99.0478515625
29 99.365234375
30 99.5849609375
31 99.0478515625
32 99.267578125
33 98.779296875
34 99.0966796875
35 98.1689453125
36 98.5595703125
37 97.0703125
38 97.265625
39 96.142578125
40 94.62890625
41 93.7744140625
42 93.26171875
43 93.212890625
44 90.9423828125
45 89.208984375
46 88.37890625
47 84.5458984375
48 91.1376953125
49 89.1357421875
50 92.4560546875
51 89.208984375
52 86.81640625
53 87.255859375
54 85.7177734375
55 84.3505859375
56 85.888671875
57 84.27734375
58 84.66796875
59 75.6103515625
60 68.212890625
61 65.1611328125
62 56.5673828125
63 48.681640625
64 48.73046875
65 37.4755859375
66 36.9384765625
67 28.759765625
68 26.953125
69 19.4091796875
70 16.50390625
71 12.841796875
72 11.8896484375
73 8.4716796875
74 6.2255859375
75 3.271484375
76 1.904296875
77 2.9296875
78 1.7578125
79 2.3193359375
80 0.78125
81 1.7333984375
82 1.8798828125
83 0.29296875
84 0.244140625
85 0.7080078125
86 0.6103515625
87 1.2939453125
88 1.806640625
89 0.9765625
90 0.830078125
91 1.416015625
92 1.07421875
93 1.07421875
94 1.025390625
95 1.1474609375
96 0.8544921875
97 1.2939453125
98 1.1474609375
99 0.4638671875
100 0.8056640625
101 1.0986328125
102 1.0009765625
103 1.0009765625
104 1.6845703125
105 0.87890625
106 0.830078125
107 1.318359375
108 0.1220703125
109 0.4638671875
110 0.390625
111 0.48828125
112 0.2685546875
113 0.1953125
114 0.3173828125
115 0.5126953125
116 0.0244140625
117 0.2197265625
118 0.341796875
119 0.2685546875
120 0.2685546875
121 0.0244140625
122 0.0732421875
123 0.0244140625
124 0.1708984375
125 0.1708984375
126 0.048828125
127 0.1220703125
128 0.0732421875
129 0.0732421875
130 0.0244140625
131 0.0244140625
132 0.0244140625
133 0
134 0.0244140625
135 0
136 0
137 0
138 0
139 0
140 0
141 0
142 0
143 0
144 0
145 0
146 0.0244140625
147 0
148 0
149 0
150 0
151 0
152 0
153 0
154 0
155 0
156 0
157 0
158 0
159 0.0244140625
160 0
161 0
162 0
163 0
164 0
165 0
166 0
167 0
168 0
169 0
170 0
171 0
172 0
173 0
174 0
175 0
176 0
177 0
178 0
179 0
180 0
181 0
182 0
183 0
184 0
185 0
186 0
187 0.0244140625
188 0
189 0
190 0
191 0
192 0
193 0.0244140625
194 0
195 0
196 0
197 0
198 0.0244140625
199 0
200 0
201 0
202 0
203 0
204 0
205 0
206 0
207 0
208 0
209 0
210 0
211 0
212 0.0244140625
213 0
214 0
215 0
216 0.0244140625
217 0
218 0
219 0
220 0
221 0
222 0
223 0
224 0
225 0
226 0
227 0
228 0
229 0
230 0
231 0
232 0
233 0
234 0
235 0
236 0
237 0
238 0
239 0
240 0
241 0
242 0
243 0
244 0
245 0
246 0
247 0
248 0
249 0
250 0
251 0
252 0.0244140625
253 0
254 0
255 0
256 0
257 0
258 0
259 0
260 0
261 0
262 0
263 0
264 0
265 0
266 0
267 0
268 0
269 0
270 0
271 0
272 0
273 0
274 0
275 0
276 0
277 0
278 0
279 0
280 0
281 0
282 0
283 0
284 0
285 0
286 0
287 0
288 0
289 0
290 0
291 0
292 0
293 0
294 0
295 0
296 0
297 0
298 0
299 0
300 0
301 0
302 0
303 0
304 0
305 0
306 0
307 0
308 0
309 0
310 0
311 0
312 0.0244140625
313 0
314 0
315 0
316 0
317 0
318 0
319 0
320 0
321 0
322 0
323 0
324 0
325 0
326 0
327 0
328 0
329 0
330 0
331 0
332 0
333 0
334 0
335 0
336 0
337 0
338 0
339 0
340 0
341 0
342 0
343 0
344 0
345 0
346 0
347 0
348 0
349 0
350 0
351 0
352 0
353 0
354 0
355 0
356 0
357 0
358 0
359 0
360 0
361 0
362 0
363 0
364 0
365 0
366 0
367 0
368 0
369 0
370 0
371 0
372 0
373 0
374 0
375 0
376 0
377 0
378 0
379 0
380 0
381 0
382 0
383 0
384 0
385 0
386 0
387 0
388 0
389 0
390 0
391 0
392 0
393 0
394 0
395 0
396 0
397 0
398 0
399 0
400 0
401 0
402 0
403 0
404 0
405 0
406 0
407 0
408 0
409 0
410 0
411 0
412 0
413 0
414 0
415 0
416 0
417 0
418 0
419 0
420 0
421 0
422 0
423 0
424 0
425 0
426 0
427 0
428 0
429 0
430 0
431 0
432 0
433 0
434 0
435 0
436 0
437 0
438 0
439 0
440 0
441 0
442 0
443 0
444 0
445 0
446 0
447 0
448 0
449 0
450 0
451 0
452 0
453 0
454 0
455 0
456 0
457 0
458 0
459 0
460 0
461 0
462 0
463 0
464 0
465 0
466 0
467 0
468 0
469 0
470 0
471 0
472 0
473 0
474 0
475 0
476 0
477 0
478 0
479 0
480 0
481 0
482 0
483 0
484 0
485 0
486 0
487 0
488 0
489 0
490 0
491 0
492 0
493 0
494 0
495 0
496 0
497 0
498 0
499 0
};
\addplot [line width=0.64pt, darkslateblue6857130]
table {%
0 100
1 100
2 100
3 100
4 100
5 100
6 100
7 100
8 99.9755859375
9 99.9755859375
10 99.9267578125
11 100
12 100
13 99.951171875
14 99.951171875
15 99.9755859375
16 99.90234375
17 99.9755859375
18 99.6337890625
19 99.51171875
20 99.6826171875
21 99.5361328125
22 99.7802734375
23 99.8291015625
24 99.7802734375
25 99.8046875
26 99.9755859375
27 99.8046875
28 98.92578125
29 99.4384765625
30 99.3896484375
31 98.92578125
32 98.5107421875
33 98.4130859375
34 97.75390625
35 96.728515625
36 97.0703125
37 96.0693359375
38 94.091796875
39 94.677734375
40 90.869140625
41 93.408203125
42 89.55078125
43 91.30859375
44 90.3564453125
45 89.9658203125
46 88.2080078125
47 88.28125
48 91.4306640625
49 90.4296875
50 88.18359375
51 89.3310546875
52 84.66796875
53 83.544921875
54 75.244140625
55 80.126953125
56 80.37109375
57 70.0927734375
58 67.822265625
59 66.064453125
60 54.736328125
61 48.8525390625
62 43.1884765625
63 40.4541015625
64 27.8076171875
65 25.6103515625
66 16.845703125
67 13.76953125
68 13.0615234375
69 9.0576171875
70 7.5927734375
71 7.8857421875
72 6.7138671875
73 5.224609375
74 4.8583984375
75 3.90625
76 0.5126953125
77 0.5126953125
78 0.29296875
79 1.5380859375
80 1.46484375
81 1.5625
82 1.07421875
83 1.1474609375
84 0.8056640625
85 0.87890625
86 1.6357421875
87 1.0498046875
88 1.0498046875
89 0.8544921875
90 0.8056640625
91 0.87890625
92 0.8544921875
93 0.537109375
94 0.439453125
95 0.6591796875
96 0.09765625
97 1.0009765625
98 0.732421875
99 0.09765625
100 1.0498046875
101 0.6103515625
102 0.927734375
103 0.8056640625
104 0.87890625
105 0.6103515625
106 0.341796875
107 0.68359375
108 0.830078125
109 0.6591796875
110 0.9033203125
111 0
112 0.390625
113 0.7080078125
114 0.5126953125
115 0.5126953125
116 0.4150390625
117 0.29296875
118 0.1953125
119 0.48828125
120 0.5859375
121 0.146484375
122 0.5859375
123 0.3173828125
124 0
125 0
126 0.0244140625
127 0
128 0
129 0.0244140625
130 0
131 0
132 0.0244140625
133 0.0244140625
134 0
135 0
136 0.0244140625
137 0
138 0
139 0
140 0
141 0
142 0
143 0
144 0
145 0
146 0
147 0
148 0
149 0.0244140625
150 0
151 0
152 0.0244140625
153 0.048828125
154 0
155 0
156 0
157 0
158 0
159 0
160 0
161 0
162 0
163 0
164 0
165 0
166 0
167 0
168 0.0244140625
169 0
170 0
171 0
172 0
173 0
174 0
175 0
176 0
177 0
178 0
179 0
180 0
181 0
182 0
183 0
184 0
185 0
186 0
187 0
188 0
189 0
190 0
191 0
192 0
193 0
194 0
195 0
196 0
197 0
198 0
199 0
200 0
201 0
202 0
203 0
204 0
205 0.0244140625
206 0
207 0
208 0
209 0
210 0
211 0
212 0
213 0
214 0.0244140625
215 0
216 0
217 0
218 0
219 0
220 0
221 0
222 0
223 0
224 0
225 0
226 0
227 0
228 0
229 0
230 0
231 0
232 0
233 0
234 0
235 0
236 0
237 0
238 0
239 0
240 0
241 0
242 0
243 0
244 0
245 0
246 0
247 0
248 0
249 0
250 0
251 0
252 0
253 0
254 0
255 0
256 0
257 0
258 0
259 0
260 0
261 0
262 0
263 0
264 0
265 0
266 0
267 0
268 0
269 0
270 0
271 0
272 0
273 0
274 0
275 0
276 0
277 0
278 0
279 0
280 0
281 0
282 0
283 0
284 0
285 0
286 0
287 0
288 0
289 0
290 0
291 0
292 0
293 0
294 0
295 0
296 0
297 0
298 0
299 0
300 0
301 0
302 0
303 0
304 0
305 0
306 0
307 0
308 0
309 0
310 0
311 0
312 0
313 0
314 0
315 0
316 0
317 0
318 0
319 0
320 0
321 0
322 0
323 0
324 0
325 0
326 0
327 0
328 0
329 0
330 0
331 0
332 0
333 0
334 0
335 0
336 0
337 0
338 0
339 0
340 0
341 0
342 0
343 0
344 0
345 0
346 0
347 0
348 0
349 0
350 0
351 0
352 0
353 0
354 0
355 0
356 0
357 0
358 0
359 0
360 0
361 0
362 0
363 0
364 0
365 0
366 0
367 0
368 0
369 0.0244140625
370 0
371 0
372 0
373 0
374 0
375 0
376 0
377 0
378 0
379 0
380 0
381 0
382 0
383 0
384 0
385 0
386 0
387 0
388 0
389 0
390 0
391 0
392 0
393 0
394 0
395 0
396 0
397 0
398 0
399 0
400 0
401 0
402 0
403 0
404 0
405 0
406 0
407 0
408 0
409 0
410 0
411 0
412 0
413 0
414 0
415 0
416 0
417 0
418 0
419 0
420 0
421 0
422 0
423 0
424 0
425 0
426 0
427 0
428 0
429 0
430 0
431 0
432 0
433 0
434 0
435 0
436 0
437 0
438 0
439 0
440 0
441 0
442 0
443 0.0244140625
444 0
445 0
446 0.0244140625
447 0
448 0
449 0
450 0
451 0
452 0.0244140625
453 0
454 0
455 0
456 0
457 0
458 0
459 0
460 0
461 0
462 0
463 0
464 0
465 0
466 0
467 0
468 0
469 0
470 0
471 0
472 0
473 0
474 0
475 0
476 0
477 0
478 0
479 0
480 0
481 0
482 0
483 0
484 0
485 0
486 0
487 0
488 0
489 0
490 0
491 0
492 0
493 0
494 0
495 0
496 0
497 0
498 0
499 0
};
\addplot [line width=0.64pt, darkslateblue5882139]
table {%
0 99.9755859375
1 100
2 100
3 99.951171875
4 99.951171875
5 99.9755859375
6 100
7 99.9755859375
8 99.951171875
9 99.951171875
10 99.9755859375
11 100
12 99.9267578125
13 99.8291015625
14 100
15 99.560546875
16 99.6826171875
17 99.8046875
18 99.9267578125
19 99.755859375
20 99.8291015625
21 99.8291015625
22 99.7314453125
23 99.70703125
24 99.8046875
25 99.6826171875
26 99.3408203125
27 98.388671875
28 98.14453125
29 98.046875
30 97.705078125
31 97.6806640625
32 96.4111328125
33 95.751953125
34 96.2646484375
35 92.7734375
36 93.017578125
37 94.3115234375
38 91.3330078125
39 89.2333984375
40 85.2783203125
41 86.669921875
42 88.28125
43 84.912109375
44 90.0634765625
45 84.912109375
46 85.15625
47 84.3017578125
48 83.203125
49 79.7119140625
50 77.392578125
51 76.5869140625
52 74.0966796875
53 72.36328125
54 67.529296875
55 63.623046875
56 53.5400390625
57 52.490234375
58 42.236328125
59 37.890625
60 25.146484375
61 22.3876953125
62 19.43359375
63 14.6728515625
64 11.9384765625
65 11.03515625
66 9.716796875
67 6.591796875
68 5.2734375
69 4.78515625
70 2.9296875
71 3.5400390625
72 2.34375
73 2.392578125
74 1.4892578125
75 2.4169921875
76 2.1728515625
77 1.025390625
78 1.220703125
79 0.5859375
80 0.146484375
81 0.8544921875
82 0.146484375
83 0.8544921875
84 1.1474609375
85 0.9765625
86 0.9765625
87 0.8544921875
88 1.025390625
89 0.8056640625
90 0.634765625
91 0.830078125
92 1.2451171875
93 0.09765625
94 0.7568359375
95 0
96 0.0244140625
97 0.048828125
98 0.1708984375
99 0.5126953125
100 0.0732421875
101 0.6103515625
102 0.6103515625
103 0.48828125
104 0.8056640625
105 0.78125
106 0.0244140625
107 0.927734375
108 0.68359375
109 0.4638671875
110 0.6103515625
111 0.7080078125
112 0.48828125
113 0.8056640625
114 0.634765625
115 0.3173828125
116 0.244140625
117 0.5859375
118 0.48828125
119 0
120 0
121 0
122 0
123 0.1953125
124 0.0732421875
125 0
126 0.0732421875
127 0.146484375
128 0.0244140625
129 0.1220703125
130 0.1220703125
131 0
132 0
133 0
134 0
135 0
136 0
137 0
138 0
139 0
140 0
141 0
142 0
143 0
144 0
145 0
146 0
147 0.0244140625
148 0
149 0
150 0
151 0
152 0
153 0
154 0
155 0
156 0
157 0
158 0
159 0
160 0
161 0
162 0.0244140625
163 0
164 0
165 0
166 0
167 0
168 0
169 0
170 0
171 0
172 0.0244140625
173 0
174 0
175 0
176 0.0244140625
177 0
178 0
179 0
180 0.0244140625
181 0
182 0
183 0.0244140625
184 0
185 0
186 0
187 0
188 0
189 0.0244140625
190 0
191 0
192 0
193 0
194 0
195 0
196 0
197 0
198 0
199 0
200 0
201 0
202 0.0244140625
203 0
204 0
205 0
206 0
207 0
208 0
209 0
210 0
211 0
212 0
213 0
214 0
215 0
216 0
217 0
218 0
219 0
220 0
221 0
222 0
223 0
224 0
225 0
226 0
227 0
228 0
229 0
230 0
231 0
232 0.0244140625
233 0
234 0.0244140625
235 0
236 0.0244140625
237 0
238 0
239 0
240 0
241 0
242 0
243 0
244 0
245 0
246 0
247 0
248 0
249 0
250 0
251 0
252 0
253 0
254 0
255 0
256 0
257 0
258 0.0244140625
259 0
260 0
261 0
262 0
263 0
264 0
265 0
266 0
267 0
268 0
269 0
270 0
271 0
272 0
273 0
274 0
275 0
276 0
277 0
278 0
279 0
280 0
281 0
282 0
283 0
284 0
285 0
286 0
287 0
288 0
289 0
290 0.0244140625
291 0
292 0
293 0
294 0
295 0
296 0
297 0
298 0
299 0
300 0.0244140625
301 0
302 0.0244140625
303 0
304 0
305 0
306 0
307 0
308 0
309 0
310 0
311 0
312 0
313 0
314 0
315 0
316 0
317 0
318 0
319 0
320 0
321 0
322 0
323 0
324 0.0244140625
325 0
326 0
327 0
328 0
329 0
330 0
331 0
332 0
333 0
334 0
335 0
336 0
337 0
338 0
339 0
340 0
341 0
342 0
343 0
344 0
345 0
346 0
347 0
348 0
349 0
350 0
351 0
352 0
353 0
354 0
355 0
356 0
357 0
358 0
359 0
360 0
361 0
362 0
363 0
364 0
365 0
366 0
367 0
368 0
369 0
370 0
371 0
372 0
373 0
374 0
375 0
376 0
377 0
378 0
379 0
380 0
381 0
382 0.0244140625
383 0
384 0
385 0
386 0
387 0
388 0
389 0
390 0
391 0
392 0
393 0
394 0
395 0
396 0
397 0
398 0
399 0
400 0
401 0.0244140625
402 0
403 0
404 0
405 0.0244140625
406 0
407 0
408 0
409 0
410 0
411 0
412 0
413 0
414 0
415 0
416 0
417 0
418 0
419 0.0244140625
420 0
421 0
422 0
423 0
424 0
425 0
426 0
427 0
428 0.0244140625
429 0
430 0
431 0
432 0
433 0
434 0
435 0
436 0
437 0
438 0
439 0
440 0.0244140625
441 0
442 0
443 0
444 0
445 0
446 0
447 0
448 0
449 0
450 0
451 0
452 0
453 0
454 0
455 0
456 0
457 0
458 0
459 0
460 0
461 0
462 0
463 0
464 0
465 0
466 0
467 0
468 0
469 0
470 0
471 0
472 0
473 0
474 0
475 0
476 0
477 0
478 0
479 0
480 0
481 0
482 0
483 0
484 0
485 0
486 0
487 0
488 0
489 0.0244140625
490 0
491 0
492 0
493 0.0244140625
494 0
495 0.0244140625
496 0
497 0
498 0
499 0
};
\addplot [line width=0.64pt, darkslateblue48103141]
table {%
0 99.9755859375
1 100
2 100
3 99.9755859375
4 100
5 99.9755859375
6 99.8779296875
7 99.951171875
8 99.951171875
9 99.8291015625
10 99.609375
11 99.560546875
12 99.9267578125
13 99.853515625
14 99.7314453125
15 99.560546875
16 99.951171875
17 99.755859375
18 99.7314453125
19 99.560546875
20 99.6337890625
21 99.365234375
22 98.9013671875
23 98.9013671875
24 99.0478515625
25 97.6318359375
26 98.0712890625
27 95.8984375
28 93.6767578125
29 94.43359375
30 92.3828125
31 92.1630859375
32 93.45703125
33 91.259765625
34 87.1826171875
35 88.916015625
36 87.109375
37 85.2783203125
38 88.2568359375
39 89.0869140625
40 84.326171875
41 85.986328125
42 88.6962890625
43 83.0810546875
44 82.8857421875
45 81.2744140625
46 74.70703125
47 75.4150390625
48 71.2646484375
49 63.7939453125
50 56.2744140625
51 53.6865234375
52 40.7470703125
53 32.0068359375
54 25.537109375
55 23.5107421875
56 17.2607421875
57 11.42578125
58 11.7431640625
59 7.7880859375
60 6.8359375
61 4.7119140625
62 5.9326171875
63 3.3203125
64 1.85546875
65 2.8564453125
66 4.0771484375
67 1.5869140625
68 1.6845703125
69 2.587890625
70 2.001953125
71 0.29296875
72 1.513671875
73 0.634765625
74 1.318359375
75 0.8056640625
76 0.048828125
77 0.1220703125
78 0.927734375
79 0.2685546875
80 0.4638671875
81 0.0732421875
82 0.1953125
83 0.5859375
84 0.5126953125
85 0.87890625
86 0
87 0.048828125
88 0.4638671875
89 0.48828125
90 0.4638671875
91 0
92 0
93 0.8544921875
94 0.8544921875
95 0.537109375
96 0.4638671875
97 0
98 0.0244140625
99 0.8056640625
100 1.171875
101 1.0009765625
102 0.09765625
103 0.8544921875
104 0.29296875
105 0.390625
106 0.732421875
107 1.025390625
108 0.0732421875
109 1.025390625
110 0.6103515625
111 0.634765625
112 0.4150390625
113 1.025390625
114 0.09765625
115 0.341796875
116 0.29296875
117 0.439453125
118 0.0244140625
119 0.048828125
120 0.1953125
121 0.3173828125
122 0.4150390625
123 0.2685546875
124 0.341796875
125 0.146484375
126 0
127 0
128 0
129 0
130 0.0244140625
131 0.048828125
132 0
133 0
134 0
135 0
136 0.0244140625
137 0.0244140625
138 0.0244140625
139 0
140 0
141 0.0244140625
142 0
143 0.0244140625
144 0
145 0
146 0
147 0
148 0
149 0
150 0
151 0.0244140625
152 0
153 0
154 0
155 0
156 0
157 0
158 0
159 0
160 0
161 0
162 0.0244140625
163 0
164 0.0244140625
165 0
166 0
167 0
168 0.0244140625
169 0
170 0
171 0
172 0
173 0
174 0
175 0
176 0
177 0
178 0
179 0
180 0
181 0
182 0
183 0
184 0
185 0
186 0.0244140625
187 0.0244140625
188 0
189 0.0244140625
190 0
191 0
192 0
193 0
194 0
195 0.0244140625
196 0
197 0
198 0
199 0
200 0.0244140625
201 0.0244140625
202 0
203 0
204 0
205 0
206 0
207 0
208 0
209 0
210 0
211 0
212 0
213 0
214 0
215 0
216 0
217 0
218 0
219 0.0244140625
220 0
221 0
222 0
223 0
224 0
225 0
226 0
227 0
228 0
229 0
230 0
231 0
232 0
233 0
234 0
235 0
236 0
237 0
238 0.0244140625
239 0
240 0
241 0
242 0
243 0
244 0.0244140625
245 0
246 0
247 0
248 0
249 0
250 0
251 0
252 0
253 0
254 0
255 0.0244140625
256 0
257 0
258 0
259 0
260 0
261 0
262 0
263 0
264 0
265 0
266 0
267 0
268 0
269 0
270 0
271 0
272 0
273 0
274 0
275 0
276 0
277 0
278 0
279 0.0244140625
280 0
281 0
282 0
283 0
284 0
285 0
286 0
287 0
288 0
289 0
290 0.0244140625
291 0
292 0
293 0
294 0
295 0
296 0
297 0
298 0
299 0
300 0
301 0
302 0
303 0
304 0
305 0
306 0
307 0
308 0
309 0
310 0
311 0
312 0
313 0
314 0
315 0
316 0
317 0
318 0.0244140625
319 0.0244140625
320 0
321 0
322 0
323 0
324 0
325 0
326 0
327 0
328 0
329 0
330 0
331 0
332 0
333 0
334 0
335 0
336 0.0244140625
337 0
338 0.0244140625
339 0
340 0.0244140625
341 0
342 0
343 0
344 0
345 0
346 0
347 0
348 0.0244140625
349 0
350 0
351 0
352 0
353 0
354 0
355 0
356 0
357 0
358 0
359 0
360 0
361 0
362 0
363 0
364 0
365 0.0244140625
366 0
367 0
368 0
369 0
370 0
371 0
372 0
373 0
374 0
375 0
376 0
377 0.0244140625
378 0
379 0
380 0
381 0
382 0
383 0
384 0
385 0
386 0
387 0
388 0
389 0
390 0
391 0
392 0.0244140625
393 0
394 0
395 0
396 0
397 0
398 0
399 0
400 0
401 0
402 0
403 0
404 0
405 0.048828125
406 0
407 0
408 0
409 0
410 0
411 0
412 0
413 0
414 0
415 0
416 0
417 0
418 0
419 0
420 0
421 0
422 0
423 0
424 0
425 0
426 0
427 0
428 0
429 0
430 0
431 0
432 0
433 0
434 0
435 0
436 0
437 0
438 0
439 0
440 0
441 0
442 0
443 0
444 0
445 0
446 0
447 0
448 0
449 0
450 0
451 0
452 0
453 0
454 0
455 0
456 0
457 0
458 0
459 0.0244140625
460 0
461 0
462 0
463 0
464 0
465 0
466 0
467 0
468 0
469 0
470 0
471 0
472 0
473 0
474 0
475 0
476 0
477 0.0244140625
478 0
479 0
480 0
481 0
482 0
483 0
484 0
485 0
486 0
487 0
488 0
489 0
490 0
491 0
492 0
493 0
494 0
495 0
496 0
497 0
498 0
499 0
};
\addplot [line width=0.64pt, teal39124142]
table {%
0 99.8046875
1 100
2 99.9755859375
3 99.951171875
4 99.8291015625
5 100
6 99.7802734375
7 99.4873046875
8 99.755859375
9 99.609375
10 99.755859375
11 99.4384765625
12 99.853515625
13 99.7802734375
14 99.9267578125
15 99.560546875
16 99.8046875
17 99.8779296875
18 99.853515625
19 99.5849609375
20 99.1455078125
21 98.8525390625
22 98.974609375
23 98.681640625
24 98.6572265625
25 97.2412109375
26 96.6552734375
27 96.0205078125
28 93.994140625
29 92.919921875
30 92.96875
31 90.869140625
32 88.2080078125
33 85.888671875
34 90.4541015625
35 84.4970703125
36 86.23046875
37 88.5009765625
38 90.0390625
39 84.326171875
40 83.0078125
41 87.98828125
42 84.1552734375
43 79.8828125
44 79.443359375
45 80.2001953125
46 72.4853515625
47 69.873046875
48 63.9404296875
49 53.076171875
50 44.8486328125
51 37.7685546875
52 30.126953125
53 25.4638671875
54 19.2626953125
55 16.40625
56 11.865234375
57 9.86328125
58 6.2255859375
59 4.7119140625
60 4.58984375
61 3.3935546875
62 3.4423828125
63 0.927734375
64 3.369140625
65 4.6875
66 1.3427734375
67 0.29296875
68 0.439453125
69 1.66015625
70 1.708984375
71 1.318359375
72 1.513671875
73 0.1708984375
74 1.07421875
75 1.0986328125
76 0.0244140625
77 0.8544921875
78 0.78125
79 0.6591796875
80 0.9521484375
81 1.1474609375
82 0.439453125
83 0.5126953125
84 0.2685546875
85 0.830078125
86 0.48828125
87 0.48828125
88 1.26953125
89 0.1220703125
90 0.537109375
91 0.5859375
92 0
93 0.68359375
94 0.1708984375
95 1.123046875
96 0.634765625
97 0.78125
98 1.1474609375
99 0.830078125
100 0.439453125
101 0.3173828125
102 0.3662109375
103 0.68359375
104 0.6591796875
105 0.7568359375
106 0.78125
107 0.6103515625
108 0.68359375
109 0.341796875
110 0.5615234375
111 0.0244140625
112 0.2197265625
113 0.5615234375
114 0.390625
115 0.1220703125
116 0.0244140625
117 0.0244140625
118 0.3173828125
119 0
120 0.2685546875
121 0.2197265625
122 0
123 0
124 0.1708984375
125 0
126 0
127 0
128 0
129 0
130 0.0244140625
131 0
132 0
133 0
134 0
135 0
136 0
137 0
138 0
139 0
140 0
141 0
142 0
143 0.0244140625
144 0
145 0
146 0
147 0.0244140625
148 0
149 0
150 0
151 0
152 0
153 0
154 0
155 0
156 0
157 0
158 0
159 0
160 0
161 0
162 0
163 0.0244140625
164 0.0244140625
165 0
166 0
167 0
168 0
169 0
170 0
171 0
172 0
173 0
174 0
175 0
176 0
177 0
178 0
179 0
180 0
181 0
182 0
183 0
184 0
185 0
186 0
187 0
188 0
189 0
190 0
191 0
192 0
193 0
194 0
195 0
196 0.0244140625
197 0
198 0
199 0
200 0
201 0
202 0
203 0.0244140625
204 0
205 0
206 0
207 0
208 0
209 0
210 0
211 0
212 0
213 0.0244140625
214 0
215 0.0244140625
216 0
217 0
218 0
219 0
220 0
221 0
222 0
223 0
224 0
225 0
226 0
227 0
228 0
229 0
230 0
231 0
232 0
233 0
234 0
235 0
236 0
237 0
238 0
239 0
240 0
241 0
242 0
243 0
244 0
245 0
246 0
247 0
248 0
249 0
250 0
251 0
252 0
253 0
254 0
255 0
256 0
257 0
258 0
259 0
260 0
261 0
262 0
263 0
264 0
265 0
266 0
267 0
268 0
269 0
270 0
271 0
272 0
273 0
274 0
275 0
276 0
277 0.0244140625
278 0
279 0
280 0
281 0
282 0
283 0
284 0.0244140625
285 0
286 0
287 0
288 0
289 0
290 0
291 0
292 0
293 0
294 0
295 0
296 0.0244140625
297 0
298 0
299 0
300 0
301 0
302 0
303 0
304 0
305 0
306 0
307 0
308 0
309 0
310 0
311 0
312 0
313 0
314 0
315 0
316 0
317 0
318 0
319 0
320 0
321 0.0244140625
322 0
323 0
324 0
325 0
326 0
327 0
328 0
329 0
330 0
331 0
332 0
333 0
334 0
335 0
336 0
337 0
338 0
339 0
340 0
341 0.0244140625
342 0
343 0
344 0
345 0
346 0
347 0
348 0
349 0
350 0
351 0
352 0
353 0
354 0
355 0
356 0
357 0
358 0
359 0
360 0
361 0
362 0
363 0
364 0.0244140625
365 0
366 0
367 0
368 0
369 0
370 0
371 0
372 0
373 0
374 0
375 0
376 0
377 0
378 0
379 0
380 0.0244140625
381 0
382 0
383 0
384 0
385 0
386 0
387 0
388 0
389 0
390 0
391 0
392 0
393 0
394 0
395 0
396 0
397 0
398 0
399 0
400 0
401 0
402 0
403 0
404 0
405 0
406 0
407 0
408 0
409 0
410 0
411 0
412 0.0244140625
413 0
414 0
415 0
416 0
417 0
418 0
419 0
420 0
421 0
422 0
423 0
424 0
425 0
426 0
427 0
428 0
429 0
430 0
431 0
432 0
433 0.0244140625
434 0
435 0
436 0
437 0
438 0
439 0
440 0
441 0
442 0
443 0
444 0
445 0
446 0
447 0
448 0
449 0
450 0
451 0
452 0
453 0
454 0
455 0
456 0.0244140625
457 0
458 0
459 0
460 0
461 0
462 0
463 0
464 0
465 0
466 0
467 0
468 0
469 0
470 0
471 0
472 0
473 0
474 0
475 0
476 0
477 0
478 0
479 0
480 0
481 0
482 0
483 0
484 0
485 0
486 0
487 0
488 0
489 0
490 0
491 0
492 0
493 0
494 0
495 0
496 0
497 0
498 0
499 0
};
\addplot [line width=0.64pt, darkcyan32144140]
table {%
0 99.9755859375
1 99.853515625
2 99.8046875
3 99.658203125
4 99.951171875
5 99.8779296875
6 99.7802734375
7 99.31640625
8 99.4140625
9 99.6337890625
10 99.755859375
11 99.31640625
12 99.5849609375
13 99.6337890625
14 99.3408203125
15 99.560546875
16 98.5107421875
17 99.5361328125
18 97.314453125
19 98.0224609375
20 98.046875
21 97.8271484375
22 96.875
23 95.8251953125
24 93.505859375
25 94.2626953125
26 91.11328125
27 91.9921875
28 86.9873046875
29 88.6474609375
30 83.740234375
31 85.0830078125
32 92.138671875
33 83.9111328125
34 86.767578125
35 91.6748046875
36 82.7880859375
37 81.4453125
38 82.5927734375
39 78.7353515625
40 70.5322265625
41 72.7294921875
42 66.845703125
43 57.3486328125
44 59.521484375
45 56.5673828125
46 47.7294921875
47 37.01171875
48 35.8154296875
49 25.830078125
50 20.8740234375
51 17.08984375
52 10.5224609375
53 9.716796875
54 6.34765625
55 5.46875
56 5.6884765625
57 3.22265625
58 6.2255859375
59 4.541015625
60 1.9775390625
61 3.6376953125
62 1.5625
63 1.7333984375
64 2.099609375
65 0.9033203125
66 0.927734375
67 2.44140625
68 1.171875
69 1.2939453125
70 1.0009765625
71 1.6357421875
72 0.927734375
73 0.5126953125
74 0.1708984375
75 1.26953125
76 0.0732421875
77 0.146484375
78 0.7080078125
79 0.048828125
80 0.78125
81 0.2197265625
82 0.87890625
83 0.7080078125
84 1.220703125
85 0.09765625
86 0.5126953125
87 0.68359375
88 0.439453125
89 0.78125
90 0.6103515625
91 0.439453125
92 0
93 0
94 0
95 0.7080078125
96 0.7080078125
97 0.634765625
98 0.48828125
99 0.3173828125
100 0.634765625
101 0.390625
102 0.7080078125
103 0.7080078125
104 0.5126953125
105 0.8544921875
106 0.7080078125
107 0.7568359375
108 0
109 1.0986328125
110 0.6591796875
111 0.4638671875
112 0.7080078125
113 0.4638671875
114 0.5126953125
115 0.4150390625
116 0.5859375
117 0.341796875
118 0
119 0.1953125
120 0.390625
121 0.341796875
122 0.29296875
123 0.4150390625
124 0
125 0
126 0
127 0
128 0
129 0.0244140625
130 0
131 0
132 0
133 0
134 0.0244140625
135 0.0244140625
136 0.0244140625
137 0
138 0
139 0
140 0.0244140625
141 0
142 0.0244140625
143 0
144 0
145 0
146 0.048828125
147 0
148 0
149 0
150 0
151 0
152 0
153 0
154 0.0244140625
155 0.0244140625
156 0
157 0
158 0
159 0.0244140625
160 0
161 0
162 0
163 0
164 0
165 0
166 0
167 0
168 0.0244140625
169 0.0244140625
170 0
171 0
172 0
173 0
174 0
175 0.0244140625
176 0.0244140625
177 0
178 0
179 0
180 0
181 0
182 0
183 0
184 0
185 0
186 0
187 0
188 0
189 0
190 0
191 0
192 0
193 0
194 0
195 0
196 0
197 0
198 0
199 0
200 0
201 0
202 0
203 0
204 0
205 0.0244140625
206 0
207 0
208 0
209 0
210 0
211 0
212 0
213 0
214 0
215 0
216 0
217 0
218 0
219 0
220 0
221 0
222 0
223 0
224 0
225 0
226 0
227 0
228 0
229 0
230 0
231 0
232 0
233 0
234 0.0244140625
235 0
236 0.0244140625
237 0
238 0
239 0
240 0
241 0
242 0
243 0
244 0
245 0
246 0
247 0
248 0
249 0
250 0
251 0
252 0
253 0
254 0
255 0
256 0
257 0
258 0
259 0
260 0
261 0
262 0
263 0.0244140625
264 0
265 0
266 0
267 0
268 0
269 0
270 0
271 0
272 0
273 0
274 0
275 0
276 0
277 0
278 0
279 0.0244140625
280 0
281 0.0244140625
282 0
283 0
284 0
285 0
286 0
287 0.0244140625
288 0
289 0
290 0
291 0
292 0
293 0
294 0
295 0
296 0
297 0
298 0
299 0
300 0
301 0
302 0
303 0
304 0
305 0
306 0
307 0
308 0
309 0
310 0
311 0
312 0
313 0
314 0
315 0
316 0
317 0
318 0
319 0
320 0
321 0
322 0
323 0
324 0
325 0
326 0
327 0
328 0
329 0
330 0
331 0
332 0
333 0
334 0
335 0
336 0
337 0
338 0
339 0
340 0
341 0
342 0
343 0
344 0
345 0
346 0
347 0
348 0
349 0
350 0
351 0
352 0
353 0
354 0
355 0
356 0.0244140625
357 0
358 0
359 0
360 0
361 0
362 0
363 0
364 0
365 0.0244140625
366 0
367 0
368 0
369 0
370 0.0244140625
371 0
372 0
373 0
374 0
375 0
376 0.0244140625
377 0
378 0
379 0
380 0
381 0
382 0
383 0
384 0
385 0
386 0
387 0
388 0
389 0
390 0
391 0
392 0
393 0
394 0
395 0
396 0
397 0
398 0
399 0
400 0
401 0
402 0
403 0
404 0
405 0
406 0
407 0
408 0
409 0
410 0
411 0
412 0
413 0
414 0
415 0
416 0
417 0
418 0
419 0
420 0
421 0
422 0
423 0
424 0
425 0
426 0
427 0
428 0
429 0
430 0
431 0
432 0
433 0
434 0
435 0
436 0
437 0.0244140625
438 0
439 0
440 0
441 0
442 0
443 0
444 0
445 0
446 0
447 0
448 0
449 0
450 0
451 0
452 0
453 0
454 0
455 0
456 0
457 0
458 0
459 0
460 0
461 0
462 0
463 0
464 0
465 0
466 0
467 0
468 0
469 0
470 0
471 0
472 0
473 0
474 0
475 0
476 0
477 0
478 0
479 0
480 0
481 0
482 0
483 0
484 0
485 0
486 0
487 0
488 0
489 0
490 0
491 0
492 0
493 0
494 0
495 0
496 0
497 0
498 0
499 0
};
\addplot [line width=0.64pt, mediumseagreen32164133]
table {%
0 99.853515625
1 99.9755859375
2 99.70703125
3 99.9267578125
4 99.4873046875
5 99.1943359375
6 99.755859375
7 99.70703125
8 99.90234375
9 99.072265625
10 99.3896484375
11 99.51171875
12 99.609375
13 98.9013671875
14 98.1689453125
15 97.6806640625
16 97.4609375
17 98.095703125
18 97.7783203125
19 96.38671875
20 91.9921875
21 95.166015625
22 93.1640625
23 93.4326171875
24 93.505859375
25 90.9423828125
26 91.30859375
27 87.0849609375
28 88.0126953125
29 86.865234375
30 90.2099609375
31 81.201171875
32 82.2998046875
33 85.791015625
34 80.6396484375
35 82.6416015625
36 80.2734375
37 78.0517578125
38 76.1474609375
39 66.2109375
40 59.1552734375
41 51.8310546875
42 41.5771484375
43 35.3759765625
44 36.81640625
45 29.833984375
46 25.48828125
47 16.8701171875
48 12.5
49 9.7412109375
50 8.4716796875
51 5.56640625
52 5.0048828125
53 4.1015625
54 1.953125
55 3.6376953125
56 3.0517578125
57 1.9287109375
58 1.85546875
59 3.7353515625
60 0.537109375
61 0.634765625
62 1.2939453125
63 1.1474609375
64 0.29296875
65 0.634765625
66 0.439453125
67 0.6103515625
68 0.4638671875
69 0.9033203125
70 2.197265625
71 1.3916015625
72 0.1220703125
73 0.8056640625
74 0.341796875
75 1.4404296875
76 0.927734375
77 0.8056640625
78 0.9033203125
79 0.439453125
80 0.5615234375
81 1.1474609375
82 0.5859375
83 0
84 1.2939453125
85 0
86 0.830078125
87 0.8056640625
88 0.439453125
89 0.3173828125
90 0.3662109375
91 0.5126953125
92 0.8544921875
93 0
94 1.025390625
95 1.123046875
96 0.8056640625
97 0.4638671875
98 0
99 0
100 0.1708984375
101 0.0244140625
102 0.2197265625
103 0.09765625
104 0.78125
105 0.439453125
106 0.5126953125
107 0.7080078125
108 1.07421875
109 0.6103515625
110 0.4638671875
111 0.5126953125
112 0.3662109375
113 0.6103515625
114 0.7568359375
115 0.830078125
116 0.2197265625
117 0.439453125
118 0.09765625
119 0.244140625
120 0.29296875
121 0.0244140625
122 0.0732421875
123 0.146484375
124 0.2685546875
125 0.048828125
126 0
127 0
128 0
129 0
130 0
131 0
132 0
133 0
134 0
135 0.0244140625
136 0
137 0.0244140625
138 0
139 0
140 0
141 0
142 0
143 0
144 0
145 0
146 0.0244140625
147 0.0244140625
148 0
149 0.0244140625
150 0
151 0
152 0
153 0.0244140625
154 0
155 0
156 0
157 0
158 0
159 0
160 0
161 0
162 0
163 0
164 0
165 0
166 0.0244140625
167 0
168 0
169 0
170 0
171 0
172 0
173 0
174 0
175 0
176 0
177 0
178 0
179 0
180 0
181 0
182 0
183 0.0244140625
184 0
185 0.0244140625
186 0.0244140625
187 0
188 0
189 0
190 0
191 0
192 0
193 0
194 0
195 0
196 0
197 0.0244140625
198 0
199 0
200 0
201 0
202 0
203 0
204 0
205 0
206 0
207 0.0244140625
208 0
209 0.0244140625
210 0
211 0
212 0
213 0
214 0
215 0
216 0
217 0.048828125
218 0.048828125
219 0
220 0
221 0
222 0
223 0
224 0
225 0
226 0
227 0
228 0
229 0
230 0
231 0
232 0
233 0.0244140625
234 0
235 0
236 0
237 0
238 0
239 0.0244140625
240 0
241 0
242 0
243 0
244 0
245 0
246 0
247 0
248 0
249 0
250 0
251 0
252 0
253 0
254 0
255 0
256 0.0244140625
257 0
258 0
259 0
260 0.0244140625
261 0.0244140625
262 0
263 0
264 0
265 0
266 0
267 0
268 0
269 0
270 0
271 0
272 0
273 0
274 0
275 0.0244140625
276 0
277 0
278 0
279 0
280 0
281 0
282 0
283 0
284 0
285 0
286 0
287 0
288 0
289 0
290 0
291 0
292 0
293 0
294 0
295 0
296 0
297 0.0244140625
298 0
299 0
300 0
301 0
302 0
303 0
304 0
305 0
306 0
307 0
308 0.0244140625
309 0
310 0
311 0
312 0
313 0
314 0
315 0
316 0
317 0
318 0
319 0.0244140625
320 0
321 0
322 0
323 0
324 0
325 0
326 0
327 0
328 0
329 0
330 0
331 0
332 0
333 0
334 0
335 0.0244140625
336 0
337 0
338 0
339 0
340 0
341 0
342 0
343 0
344 0
345 0
346 0
347 0
348 0
349 0
350 0
351 0
352 0
353 0
354 0
355 0
356 0
357 0
358 0
359 0
360 0
361 0
362 0
363 0
364 0
365 0
366 0
367 0
368 0
369 0
370 0
371 0
372 0
373 0
374 0
375 0
376 0
377 0
378 0
379 0
380 0
381 0
382 0
383 0
384 0
385 0
386 0
387 0
388 0
389 0
390 0
391 0
392 0
393 0
394 0
395 0
396 0
397 0
398 0
399 0
400 0
401 0
402 0
403 0
404 0
405 0
406 0
407 0
408 0
409 0
410 0
411 0
412 0
413 0
414 0
415 0
416 0
417 0
418 0
419 0
420 0
421 0.0244140625
422 0
423 0
424 0
425 0
426 0
427 0
428 0
429 0
430 0
431 0
432 0
433 0
434 0
435 0
436 0
437 0
438 0
439 0
440 0
441 0
442 0
443 0
444 0
445 0
446 0
447 0
448 0
449 0
450 0
451 0
452 0
453 0
454 0
455 0
456 0
457 0
458 0
459 0
460 0
461 0
462 0
463 0
464 0
465 0
466 0
467 0
468 0
469 0
470 0
471 0
472 0
473 0
474 0
475 0
476 0
477 0
478 0
479 0
480 0
481 0
482 0
483 0
484 0
485 0
486 0
487 0
488 0
489 0
490 0
491 0
492 0
493 0
494 0
495 0
496 0
497 0
498 0
499 0
};
\addplot [line width=0.64pt, mediumseagreen54184119]
table {%
0 99.9267578125
1 99.267578125
2 99.9267578125
3 99.658203125
4 99.1455078125
5 99.8046875
6 98.5595703125
7 99.51171875
8 99.6337890625
9 98.6572265625
10 99.1455078125
11 99.072265625
12 98.9013671875
13 97.7783203125
14 95.5322265625
15 97.216796875
16 93.5302734375
17 93.4326171875
18 90.185546875
19 88.7939453125
20 90.0634765625
21 88.818359375
22 93.6767578125
23 90.966796875
24 86.6455078125
25 87.6708984375
26 88.3056640625
27 89.5263671875
28 81.1767578125
29 82.470703125
30 80.2734375
31 85.7177734375
32 84.619140625
33 77.5390625
34 77.1240234375
35 68.0908203125
36 63.37890625
37 56.9580078125
38 43.2373046875
39 36.1083984375
40 31.0302734375
41 20.849609375
42 15.3564453125
43 12.890625
44 9.9365234375
45 8.740234375
46 5.46875
47 7.2265625
48 2.4658203125
49 1.7333984375
50 4.3212890625
51 1.0009765625
52 0.5126953125
53 1.9287109375
54 0.2685546875
55 0.0244140625
56 0.341796875
57 0.146484375
58 0.2685546875
59 0.390625
60 0.2197265625
61 2.24609375
62 1.46484375
63 1.8798828125
64 1.5380859375
65 1.85546875
66 0.2685546875
67 0.634765625
68 1.5869140625
69 1.0009765625
70 1.220703125
71 0.732421875
72 1.0986328125
73 0.5615234375
74 0.09765625
75 0.9765625
76 0.0244140625
77 1.171875
78 0.048828125
79 0.048828125
80 0
81 0
82 0.537109375
83 0.048828125
84 0.0244140625
85 0
86 0
87 0
88 0
89 0.0244140625
90 0
91 0
92 0.0244140625
93 0.9033203125
94 0
95 1.123046875
96 0.048828125
97 0.341796875
98 0
99 0.5859375
100 0.9521484375
101 0.48828125
102 0.732421875
103 0.9765625
104 0.6591796875
105 0
106 0.87890625
107 0.537109375
108 1.025390625
109 0.48828125
110 0.390625
111 0.4638671875
112 0.87890625
113 0.439453125
114 0.78125
115 0.5615234375
116 0.5126953125
117 0.3173828125
118 0.2197265625
119 0.4150390625
120 0.048828125
121 0.341796875
122 0
123 0.29296875
124 0.146484375
125 0
126 0
127 0
128 0
129 0
130 0
131 0
132 0
133 0
134 0
135 0
136 0
137 0
138 0
139 0
140 0
141 0
142 0.0244140625
143 0
144 0
145 0
146 0.0244140625
147 0
148 0
149 0
150 0
151 0
152 0
153 0.0244140625
154 0
155 0
156 0
157 0
158 0
159 0
160 0
161 0.0244140625
162 0
163 0
164 0
165 0
166 0
167 0
168 0
169 0
170 0
171 0.0244140625
172 0
173 0
174 0
175 0
176 0.048828125
177 0
178 0
179 0
180 0
181 0
182 0
183 0
184 0
185 0
186 0.0244140625
187 0
188 0.0244140625
189 0
190 0
191 0
192 0
193 0.0244140625
194 0
195 0
196 0
197 0
198 0
199 0
200 0.0244140625
201 0
202 0
203 0
204 0
205 0
206 0
207 0
208 0.0244140625
209 0
210 0
211 0
212 0
213 0
214 0
215 0
216 0
217 0
218 0
219 0.0244140625
220 0
221 0
222 0
223 0
224 0.0244140625
225 0
226 0
227 0.0244140625
228 0
229 0
230 0.0244140625
231 0
232 0
233 0
234 0
235 0
236 0
237 0
238 0.048828125
239 0
240 0
241 0
242 0.0244140625
243 0
244 0
245 0.048828125
246 0
247 0
248 0
249 0.0244140625
250 0
251 0.0244140625
252 0
253 0
254 0
255 0
256 0
257 0
258 0
259 0
260 0
261 0
262 0
263 0.0244140625
264 0.0244140625
265 0
266 0.0244140625
267 0
268 0
269 0
270 0
271 0
272 0
273 0
274 0
275 0
276 0
277 0
278 0
279 0
280 0
281 0
282 0.0244140625
283 0
284 0
285 0
286 0
287 0
288 0
289 0
290 0
291 0
292 0
293 0
294 0
295 0
296 0
297 0
298 0
299 0
300 0
301 0
302 0
303 0
304 0
305 0
306 0
307 0
308 0
309 0.0244140625
310 0
311 0
312 0
313 0
314 0
315 0
316 0
317 0
318 0
319 0
320 0
321 0
322 0
323 0
324 0
325 0
326 0
327 0.0244140625
328 0
329 0
330 0
331 0
332 0
333 0
334 0
335 0
336 0.0244140625
337 0
338 0
339 0.0244140625
340 0
341 0
342 0
343 0
344 0
345 0
346 0
347 0
348 0
349 0
350 0
351 0
352 0
353 0
354 0
355 0
356 0
357 0
358 0
359 0
360 0
361 0
362 0
363 0
364 0
365 0
366 0
367 0
368 0
369 0
370 0
371 0
372 0
373 0
374 0
375 0
376 0
377 0
378 0
379 0
380 0
381 0
382 0.0244140625
383 0.0244140625
384 0
385 0
386 0
387 0
388 0
389 0
390 0
391 0
392 0
393 0
394 0
395 0
396 0
397 0
398 0
399 0
400 0
401 0
402 0
403 0
404 0.0244140625
405 0
406 0
407 0
408 0
409 0
410 0
411 0
412 0
413 0
414 0
415 0
416 0
417 0
418 0
419 0
420 0
421 0
422 0
423 0
424 0
425 0
426 0
427 0
428 0
429 0
430 0
431 0
432 0
433 0
434 0
435 0
436 0
437 0
438 0
439 0
440 0
441 0
442 0
443 0
444 0
445 0
446 0
447 0
448 0
449 0
450 0
451 0
452 0
453 0.0244140625
454 0
455 0
456 0
457 0
458 0
459 0
460 0
461 0
462 0
463 0
464 0
465 0
466 0
467 0
468 0
469 0
470 0
471 0
472 0
473 0
474 0.0244140625
475 0.0244140625
476 0
477 0
478 0
479 0
480 0
481 0
482 0
483 0
484 0
485 0
486 0
487 0
488 0
489 0
490 0
491 0
492 0
493 0
494 0
495 0.0244140625
496 0
497 0
498 0
499 0
};
\addplot [line width=0.64pt, mediumseagreen9420197]
table {%
0 99.90234375
1 99.4873046875
2 99.5849609375
3 99.462890625
4 99.4873046875
5 99.462890625
6 99.8779296875
7 99.072265625
8 97.998046875
9 98.8037109375
10 98.193359375
11 97.7294921875
12 96.09375
13 96.5576171875
14 95.1904296875
15 97.55859375
16 94.8486328125
17 94.82421875
18 87.6708984375
19 96.5087890625
20 92.626953125
21 93.6279296875
22 88.2568359375
23 89.4775390625
24 88.57421875
25 89.697265625
26 89.6484375
27 86.9873046875
28 76.904296875
29 79.8583984375
30 83.251953125
31 80.46875
32 75.244140625
33 71.5087890625
34 70.8984375
35 61.669921875
36 53.466796875
37 46.2646484375
38 37.2314453125
39 31.0546875
40 22.2900390625
41 15.625
42 12.1826171875
43 8.4716796875
44 9.228515625
45 7.8857421875
46 6.884765625
47 4.5166015625
48 2.1484375
49 1.7333984375
50 1.6357421875
51 2.490234375
52 1.4404296875
53 2.0751953125
54 0.2197265625
55 2.24609375
56 0.146484375
57 0.09765625
58 0.1220703125
59 0.1220703125
60 0.0732421875
61 0.634765625
62 0.048828125
63 0
64 1.1474609375
65 1.611328125
66 1.85546875
67 0.1708984375
68 0.048828125
69 1.8310546875
70 0.341796875
71 0.390625
72 0.8056640625
73 0.732421875
74 0.146484375
75 0.732421875
76 0.048828125
77 0.9521484375
78 0.68359375
79 0.7568359375
80 0.732421875
81 1.0009765625
82 0.8544921875
83 0
84 0
85 1.1962890625
86 0.0244140625
87 0.09765625
88 0.0244140625
89 0.6103515625
90 1.025390625
91 0.5126953125
92 1.46484375
93 0.3173828125
94 1.0009765625
95 0.439453125
96 0.9033203125
97 0
98 0.7568359375
99 0.6591796875
100 0.3662109375
101 0.5126953125
102 0.3662109375
103 0.634765625
104 0.7568359375
105 0.537109375
106 1.025390625
107 0.9521484375
108 0.8544921875
109 0.4150390625
110 1.0009765625
111 0.87890625
112 0.1953125
113 0.4150390625
114 0.390625
115 0.537109375
116 0.439453125
117 0.48828125
118 0.439453125
119 0
120 0.3173828125
121 0.6591796875
122 0.2197265625
123 0.3173828125
124 0
125 0.048828125
126 0
127 0
128 0
129 0
130 0
131 0
132 0
133 0
134 0
135 0.048828125
136 0.0244140625
137 0.0244140625
138 0.048828125
139 0
140 0
141 0
142 0
143 0
144 0
145 0.0244140625
146 0
147 0
148 0
149 0
150 0
151 0
152 0
153 0
154 0
155 0
156 0
157 0
158 0
159 0
160 0
161 0
162 0
163 0
164 0
165 0
166 0
167 0
168 0
169 0
170 0
171 0
172 0
173 0
174 0
175 0
176 0
177 0
178 0
179 0.0244140625
180 0
181 0
182 0
183 0
184 0
185 0
186 0.0244140625
187 0
188 0
189 0
190 0
191 0.0244140625
192 0
193 0
194 0.0244140625
195 0
196 0
197 0
198 0
199 0
200 0
201 0
202 0
203 0
204 0
205 0
206 0
207 0.0244140625
208 0
209 0
210 0
211 0
212 0
213 0
214 0.0244140625
215 0
216 0
217 0
218 0
219 0
220 0
221 0
222 0
223 0
224 0
225 0
226 0
227 0
228 0
229 0
230 0
231 0
232 0
233 0
234 0
235 0.0244140625
236 0.0244140625
237 0
238 0
239 0
240 0.0244140625
241 0
242 0
243 0
244 0
245 0
246 0
247 0
248 0
249 0
250 0
251 0
252 0.0244140625
253 0
254 0
255 0
256 0.0244140625
257 0
258 0
259 0
260 0
261 0
262 0
263 0.0244140625
264 0
265 0.0244140625
266 0
267 0
268 0
269 0
270 0
271 0.0244140625
272 0
273 0
274 0
275 0
276 0
277 0
278 0
279 0
280 0
281 0
282 0
283 0
284 0
285 0
286 0
287 0
288 0
289 0
290 0
291 0
292 0
293 0
294 0
295 0
296 0
297 0
298 0.0244140625
299 0
300 0.0244140625
301 0
302 0.0244140625
303 0
304 0
305 0
306 0
307 0
308 0.0244140625
309 0
310 0
311 0
312 0
313 0.0244140625
314 0
315 0
316 0
317 0
318 0
319 0
320 0
321 0
322 0
323 0
324 0
325 0
326 0
327 0.0244140625
328 0
329 0
330 0
331 0
332 0
333 0.0244140625
334 0
335 0
336 0
337 0
338 0
339 0
340 0
341 0
342 0
343 0
344 0
345 0
346 0
347 0
348 0
349 0
350 0
351 0
352 0
353 0
354 0
355 0
356 0
357 0
358 0
359 0
360 0
361 0
362 0
363 0
364 0
365 0
366 0
367 0
368 0
369 0
370 0
371 0
372 0
373 0
374 0
375 0
376 0
377 0
378 0
379 0
380 0
381 0
382 0
383 0
384 0
385 0
386 0
387 0
388 0
389 0
390 0
391 0
392 0
393 0
394 0.0244140625
395 0
396 0
397 0.0244140625
398 0
399 0
400 0.0244140625
401 0
402 0
403 0
404 0
405 0
406 0
407 0
408 0
409 0
410 0
411 0
412 0
413 0
414 0
415 0
416 0
417 0.0244140625
418 0
419 0
420 0
421 0
422 0
423 0
424 0
425 0
426 0
427 0
428 0
429 0
430 0
431 0
432 0
433 0
434 0
435 0
436 0
437 0
438 0
439 0
440 0
441 0
442 0
443 0
444 0
445 0
446 0
447 0
448 0
449 0
450 0
451 0.0244140625
452 0
453 0
454 0
455 0
456 0
457 0
458 0
459 0.0244140625
460 0
461 0
462 0
463 0
464 0
465 0
466 0
467 0.0244140625
468 0
469 0
470 0.0244140625
471 0
472 0.0244140625
473 0
474 0
475 0
476 0
477 0
478 0
479 0
480 0
481 0
482 0.0244140625
483 0
484 0
485 0
486 0
487 0
488 0
489 0
490 0
491 0.0244140625
492 0
493 0
494 0
495 0
496 0
497 0
498 0
499 0
};
\addplot [line width=0.64pt, yellowgreen14621565]
table {%
0 98.9501953125
1 98.8525390625
2 98.5595703125
3 96.8017578125
4 97.4365234375
5 96.044921875
6 96.77734375
7 97.119140625
8 93.9453125
9 92.4560546875
10 92.87109375
11 91.2841796875
12 86.1083984375
13 90.7958984375
14 89.013671875
15 85.7421875
16 87.158203125
17 91.259765625
18 88.5009765625
19 89.3310546875
20 82.275390625
21 80.4443359375
22 78.662109375
23 76.953125
24 72.021484375
25 62.4755859375
26 60.44921875
27 50.634765625
28 39.6484375
29 33.6181640625
30 24.9755859375
31 21.3134765625
32 15.72265625
33 11.81640625
34 8.1298828125
35 5.0048828125
36 4.0283203125
37 5.5419921875
38 2.490234375
39 1.0009765625
40 1.025390625
41 1.8310546875
42 0.4638671875
43 0.2197265625
44 0.927734375
45 2.3681640625
46 0.29296875
47 0.9765625
48 1.806640625
49 1.4892578125
50 0.1953125
51 1.66015625
52 1.26953125
53 1.0009765625
54 1.3916015625
55 0.244140625
56 2.1240234375
57 2.001953125
58 1.4404296875
59 0.390625
60 0.146484375
61 0.7568359375
62 1.025390625
63 0.5126953125
64 0.4638671875
65 1.025390625
66 1.0009765625
67 0.439453125
68 0.5126953125
69 1.3916015625
70 0.732421875
71 0.78125
72 0.048828125
73 0.537109375
74 0.5859375
75 0.0244140625
76 0.0732421875
77 0.9521484375
78 1.0009765625
79 0.5859375
80 0.4638671875
81 0.634765625
82 1.025390625
83 0.0732421875
84 0.09765625
85 0.0244140625
86 0.0244140625
87 0.537109375
88 0
89 0.5615234375
90 0.6591796875
91 0.78125
92 1.26953125
93 0.830078125
94 0.68359375
95 0.87890625
96 0.87890625
97 0.634765625
98 0.5859375
99 0.4150390625
100 0.9033203125
101 1.0498046875
102 0.6103515625
103 1.123046875
104 1.123046875
105 0.4638671875
106 0.4638671875
107 0.5615234375
108 0.1708984375
109 1.0986328125
110 0.830078125
111 0.4638671875
112 0.9033203125
113 0.3173828125
114 0.5126953125
115 0.537109375
116 0.634765625
117 0.341796875
118 0.390625
119 0.341796875
120 0.390625
121 0.244140625
122 0.5126953125
123 0.244140625
124 0.048828125
125 0.146484375
126 0
127 0.1220703125
128 0
129 0
130 0.0244140625
131 0
132 0.0732421875
133 0
134 0
135 0
136 0.0244140625
137 0
138 0.0244140625
139 0.0244140625
140 0
141 0.0244140625
142 0
143 0
144 0
145 0
146 0
147 0
148 0
149 0
150 0
151 0
152 0
153 0
154 0
155 0
156 0
157 0
158 0
159 0
160 0
161 0
162 0
163 0
164 0
165 0
166 0.0244140625
167 0
168 0
169 0
170 0
171 0
172 0
173 0
174 0
175 0
176 0
177 0
178 0
179 0
180 0
181 0
182 0
183 0
184 0
185 0
186 0
187 0
188 0
189 0
190 0
191 0
192 0
193 0
194 0
195 0
196 0
197 0
198 0
199 0
200 0
201 0
202 0
203 0
204 0
205 0
206 0
207 0
208 0
209 0
210 0
211 0
212 0
213 0
214 0
215 0
216 0
217 0
218 0
219 0
220 0
221 0.0244140625
222 0
223 0
224 0
225 0
226 0
227 0
228 0
229 0
230 0
231 0
232 0
233 0
234 0
235 0
236 0
237 0
238 0
239 0
240 0
241 0
242 0
243 0
244 0
245 0
246 0
247 0
248 0
249 0
250 0
251 0
252 0
253 0
254 0
255 0
256 0
257 0
258 0.0244140625
259 0
260 0
261 0
262 0
263 0.0244140625
264 0
265 0
266 0
267 0
268 0
269 0
270 0
271 0
272 0
273 0
274 0
275 0
276 0
277 0
278 0
279 0
280 0
281 0
282 0
283 0
284 0
285 0
286 0
287 0
288 0
289 0
290 0
291 0
292 0
293 0
294 0
295 0
296 0
297 0
298 0
299 0
300 0.0244140625
301 0
302 0
303 0
304 0
305 0
306 0
307 0
308 0
309 0
310 0
311 0
312 0
313 0
314 0
315 0
316 0
317 0
318 0
319 0
320 0
321 0
322 0
323 0
324 0
325 0
326 0
327 0
328 0
329 0
330 0.0244140625
331 0
332 0
333 0
334 0
335 0
336 0
337 0
338 0
339 0
340 0
341 0
342 0
343 0
344 0
345 0
346 0
347 0
348 0
349 0
350 0
351 0
352 0
353 0
354 0
355 0
356 0.0244140625
357 0.0244140625
358 0.0244140625
359 0
360 0
361 0
362 0
363 0
364 0
365 0
366 0.0244140625
367 0
368 0
369 0
370 0.0244140625
371 0.048828125
372 0
373 0
374 0
375 0
376 0
377 0
378 0
379 0
380 0.0244140625
381 0
382 0
383 0
384 0
385 0.0244140625
386 0
387 0
388 0
389 0
390 0
391 0
392 0
393 0.0244140625
394 0
395 0
396 0
397 0
398 0
399 0
400 0
401 0
402 0
403 0
404 0
405 0
406 0
407 0
408 0
409 0
410 0
411 0
412 0
413 0
414 0
415 0
416 0.0244140625
417 0
418 0
419 0
420 0
421 0
422 0
423 0.0244140625
424 0
425 0
426 0
427 0
428 0
429 0
430 0
431 0
432 0
433 0
434 0
435 0
436 0
437 0.0244140625
438 0
439 0
440 0
441 0
442 0
443 0
444 0
445 0
446 0
447 0
448 0
449 0
450 0
451 0
452 0
453 0
454 0
455 0
456 0
457 0
458 0.0244140625
459 0
460 0
461 0
462 0
463 0
464 0
465 0
466 0
467 0
468 0
469 0
470 0
471 0
472 0
473 0
474 0
475 0
476 0
477 0
478 0
479 0
480 0
481 0
482 0
483 0
484 0
485 0
486 0
487 0
488 0
489 0
490 0
491 0
492 0
493 0
494 0
495 0
496 0
497 0
498 0
499 0
};
\addplot [line width=0.64pt, greenyellow20222430]
table {%
0 98.9013671875
1 97.509765625
2 96.337890625
3 97.509765625
4 97.6806640625
5 93.8720703125
6 94.53125
7 90.5029296875
8 89.306640625
9 89.5751953125
10 89.306640625
11 92.67578125
12 93.4814453125
13 89.6240234375
14 84.7900390625
15 83.6181640625
16 84.5703125
17 85.888671875
18 78.076171875
19 76.318359375
20 71.09375
21 70.947265625
22 62.79296875
23 55.908203125
24 43.994140625
25 37.6220703125
26 31.7138671875
27 25.9765625
28 15.0634765625
29 12.3779296875
30 7.763671875
31 6.396484375
32 4.0283203125
33 1.953125
34 1.5625
35 0.8056640625
36 0.8544921875
37 0.537109375
38 1.171875
39 1.220703125
40 1.6357421875
41 2.2705078125
42 1.3671875
43 2.1728515625
44 0.244140625
45 1.4892578125
46 1.9775390625
47 0.1708984375
48 1.611328125
49 0.4150390625
50 0.4150390625
51 0.048828125
52 0.4638671875
53 0.390625
54 1.3916015625
55 1.025390625
56 1.6845703125
57 0.9033203125
58 0.3662109375
59 0.244140625
60 0.0732421875
61 0
62 1.3671875
63 0.1220703125
64 0.732421875
65 0.3662109375
66 0.6591796875
67 0.830078125
68 0.390625
69 0.0244140625
70 0.09765625
71 0.244140625
72 0.0244140625
73 0.78125
74 0.3173828125
75 0.537109375
76 0
77 0.8544921875
78 0.48828125
79 0
80 0.390625
81 0.439453125
82 0.2197265625
83 0.87890625
84 0.6103515625
85 0.0732421875
86 0.537109375
87 0.7568359375
88 0.48828125
89 0.2197265625
90 1.3916015625
91 0.48828125
92 0.8056640625
93 1.1474609375
94 0.87890625
95 0.9033203125
96 0.634765625
97 0.68359375
98 0.5615234375
99 0.634765625
100 0.2685546875
101 0.732421875
102 0.4638671875
103 0.87890625
104 0.830078125
105 0.4150390625
106 0.5615234375
107 0.8056640625
108 0.0244140625
109 0.0244140625
110 0.8056640625
111 0.3662109375
112 0.146484375
113 0
114 0.439453125
115 0
116 0.1953125
117 0.4150390625
118 0.3173828125
119 0
120 0
121 0
122 0
123 0.09765625
124 0.2197265625
125 0
126 0.0732421875
127 0
128 0
129 0
130 0
131 0
132 0
133 0
134 0
135 0.0244140625
136 0.048828125
137 0
138 0
139 0.0244140625
140 0.0244140625
141 0
142 0
143 0
144 0
145 0
146 0
147 0
148 0
149 0
150 0
151 0
152 0
153 0
154 0
155 0
156 0.0244140625
157 0
158 0.0244140625
159 0
160 0.0244140625
161 0
162 0.0244140625
163 0
164 0
165 0
166 0
167 0
168 0
169 0
170 0
171 0
172 0
173 0
174 0
175 0
176 0
177 0
178 0.0244140625
179 0
180 0
181 0
182 0
183 0
184 0
185 0
186 0
187 0
188 0
189 0
190 0
191 0
192 0
193 0
194 0
195 0
196 0
197 0
198 0
199 0
200 0
201 0
202 0
203 0
204 0
205 0
206 0
207 0
208 0
209 0
210 0
211 0
212 0
213 0
214 0
215 0
216 0
217 0
218 0
219 0
220 0
221 0
222 0
223 0
224 0
225 0
226 0
227 0
228 0
229 0
230 0
231 0
232 0
233 0
234 0
235 0
236 0
237 0
238 0
239 0
240 0
241 0
242 0
243 0
244 0
245 0
246 0
247 0
248 0
249 0
250 0.0244140625
251 0
252 0
253 0.048828125
254 0
255 0
256 0
257 0
258 0
259 0
260 0
261 0.0244140625
262 0
263 0
264 0
265 0
266 0
267 0
268 0.0244140625
269 0
270 0.0244140625
271 0
272 0.0244140625
273 0
274 0
275 0
276 0.0244140625
277 0
278 0
279 0
280 0
281 0
282 0
283 0
284 0
285 0
286 0.0244140625
287 0
288 0
289 0
290 0
291 0
292 0.0244140625
293 0
294 0.0244140625
295 0
296 0
297 0
298 0
299 0
300 0
301 0
302 0
303 0
304 0
305 0
306 0
307 0
308 0
309 0
310 0
311 0
312 0
313 0
314 0
315 0
316 0
317 0
318 0
319 0
320 0
321 0
322 0
323 0
324 0
325 0
326 0.048828125
327 0
328 0
329 0.0244140625
330 0
331 0
332 0
333 0
334 0
335 0
336 0
337 0
338 0
339 0
340 0
341 0
342 0
343 0
344 0
345 0
346 0
347 0
348 0
349 0
350 0
351 0
352 0
353 0
354 0
355 0
356 0
357 0
358 0.0244140625
359 0
360 0
361 0
362 0
363 0
364 0.0244140625
365 0
366 0
367 0
368 0
369 0
370 0
371 0
372 0
373 0
374 0
375 0
376 0
377 0
378 0
379 0
380 0
381 0
382 0
383 0
384 0
385 0
386 0
387 0
388 0
389 0
390 0
391 0
392 0
393 0
394 0
395 0
396 0
397 0
398 0
399 0
400 0
401 0
402 0
403 0
404 0
405 0
406 0
407 0
408 0
409 0
410 0
411 0
412 0
413 0
414 0
415 0.0244140625
416 0
417 0
418 0
419 0
420 0
421 0
422 0
423 0
424 0
425 0
426 0
427 0
428 0
429 0
430 0
431 0
432 0
433 0
434 0
435 0
436 0
437 0.0244140625
438 0
439 0
440 0
441 0
442 0
443 0
444 0
445 0
446 0
447 0
448 0
449 0
450 0
451 0
452 0
453 0
454 0
455 0.0244140625
456 0
457 0
458 0
459 0
460 0
461 0
462 0
463 0
464 0
465 0
466 0
467 0
468 0
469 0
470 0
471 0
472 0
473 0
474 0
475 0
476 0
477 0
478 0
479 0
480 0.0244140625
481 0
482 0
483 0
484 0
485 0
486 0
487 0
488 0
489 0
490 0
491 0
492 0
493 0
494 0
495 0
496 0
497 0
498 0
499 0
};

\end{groupplot}

\end{tikzpicture}

  \caption{Win rate of \tprobe with the probe address in different LLC slices of an Intel \labfourtytwoCPU ($n = 4096$ per sample). The threads run on cores 0 and 1, which are SMT siblings.}
  \label{fig:llc-slices}
\end{figure}
The LLCs of multicore Intel CPUs are segmented into \textit{slices}, each typically grouped with one physical CPU core~\cite{Irazoqui2015slice,Gerlach2024Hash,Rainer2025slice,Morgan2025slice}.
Slices are interconnected in a ring- or mesh-like routing structure, and LLC access latency depends on how far a slice is from the core issuing the operation.
A complex microarchitectural hash function over the physical address bits determines which slice an address uses.
We show that with multiple aggregated measurements, \AttackName can determine an address's LLC slice.
For this experiment, we use the Intel \labfourtytwoCPU, which features \num{12} LLC slices on a bidirectional ring interconnect.
The hash function on this CPU is known~\cite{Rainer2025slice}, giving us reliable ground truth.
To reliably evict addresses to the LLC, we use the eviction strategy outlined in \Cref{sec:reliable-cache-eviction}.
We first measure \tprobe's win rate at different padding levels, with the probe address in different LLC slices and the threads on SMT siblings.
The results of this experiment are shown in \Cref{fig:llc-slices}.
We observe that the point at which the win rate drops below \SI{50}{\percent} is different for each LLC slice, with the relative differences varying from \num{45} \texttt{add}-cycles for slices \num{0} and \num{11}  to just \num{1} \texttt{add}-cycle for slices \num{8} and \num{9}.
However, we also observe that no padding level creates a scenario where, e.g., an address from slice \num{8} always causes \tprobe to win and an address from slice \num{9} always causes it to lose.
To determine the cache slice, we therefore use the location of the win rate's \SI{50}{\percent} threshold as a stable marker.
Hence, we must take multiple measurements per test address with different padding levels.
We measure \AttackName's performance for determining an address's cache slice in a separate experiment.
It consists of two steps: A \emph{calibration phase}, which determines the \SI{50}{\percent} thresholds of the previous experiment using addresses with known cache slices, and an \emph{evaluation phase}, which determines the cache slice of \num{1000} random addresses and compares it to the ground truth.
In the calibration phase, we determine the \SI{50}{\percent} thresholds with the threads pinned to SMT sibling cores \num{0} and \num{1} and cores \num{14} and \num{15}, giving us measurements from two points on the ring.
The evaluation phase mirrors this, determining the \SI{50}{\percent} threshold of a test address in both configurations, using \num{512} measurements per padding level.
To determine the slice, we treat the calibration thresholds and the two test-address thresholds as points in a two-dimensional space, and choose the slice whose calibration point is closest to the observed test-address thresholds.
Using this technique, we observe a slice detection accuracy of \SI{96.9}{\percent} with \num{1000} random addresses.
This means that \AttackName can distinguish differences on the order of individual cycles in scenarios where multiple measurements are permitted.
The uncertain range for SMT siblings of this CPU in \Cref{subsec:delta-sync} is \num{31} \texttt{add}-cycles wide, confirming that multiple measurements can further improve the temporal resolution.

\subsection{Beyond Cache-Access Latencies}\label{subsec:other-latency-types}

While this paper focuses mostly on cache access latencies, \AttackName can, in principle, distinguish any two microarchitectural timing differences.
We also show it can profile other stateful components, such as the TLB and parts of the execution engine.

\paragrabf{AVX Wakeup Throttling}
On many Intel CPUs, 256-bit AVX instructions executed after a prolonged period of non-AVX usage are subject to a transient performance penalty.
In this ``cold'' state, latencies of AVX instructions are significantly higher~\cite{Schwarz2019netspectre}.
Haj-Yahya \etal show that this is due to various current management mechanisms that throttle SIMD instruction throughput while the processor adjusts voltage and frequency settings~\cite{haj2021ichannels}.
Padded \AttackName can detect this penalty in AVX-512 instructions.
For this experiment, we replace the memory access to the probe address in \tprobe with an AVX-512 fused multiply-add instruction.
Specifically, we add a single \texttt{vfnmsub231pd} instruction with two \texttt{vmov} instructions to transfer input and output between SIMD and GP registers.
We also remove the reference memory access in \tref, relying entirely on the padding in \tprobe to separate the two cases.
For the fast case, we zero all AVX registers and then pre-warm the AVX-512 units by spinning on a series of AVX-512 instructions for \SI{5}{\micro\second} before measurement.
For the slow case, we spin on \texttt{rdtsc} for \SI{1}{\milli\second}, not executing any AVX instructions.
We run this experiment on an Intel \labeightynineCPU with the threads pinned to cores 0 and 1 (i.e., separate physical cores), taking \num{512} measurement samples per case and padding level.
Starting at a padding level of \num{159} \texttt{add}-cycles, we observe a golden range of \num{222} \texttt{add}-cycles in which the combined $F_1$ score is above \num{0.98}.
In all but \num{5} outlier cases in this range, the $F_1$ score is $> 0.995$.
This shows that \AttackName can reliably detect this timing difference.

\paragrabf{FP Operations with Subnormal Operands}
Floating-point standards such as IEEE-754~\cite{ieee7542008} define a special number class known as subnormal (or sometimes called denormal) numbers.
In the binary interchange formats, a subnormal is encoded with an all-zero biased exponent field and a non-zero trailing significand, allowing the representation of values very close to zero.
This fills an otherwise unused part of the encoding space but adds complexity when implementing IEEE-754-compliant floating-point arithmetic in hardware.
Operations involving subnormal numbers often incur a performance penalty~\cite{Andrysco2015subnormal,kohlbrenner2017effectiveness}, with Intel CPUs typically relying on microcode assists for their implementation~\cite{Intel2023optimization}.
To measure this with \AttackName, we replace the memory access to the probe address in \tprobe with an AVX2 floating-point division, and let it operate on normal numbers for the fast case and subnormal numbers for the slow case.
As in the previous experiment, we remove the reference memory access in \tref.
We run this experiment on SMT sibling cores of an Intel \labeightynineCPU, taking \num{4096} samples per configuration and padding level.
We observe an $F_1$ score of \num{0.9991} at padding level $0$, with the golden range extending to \num{31} \texttt{add}-cycles.
Hence, this timing difference can be reliably detected even with unpadded \AttackName.

\paragrabf{TLB Residency}
\AttackName can detect whether a CPU's TLB caches a translation for a given virtual address.
We verify this with an experiment on an Intel \labfourtytwoCPU.
In the fast case, we access the probe address before measurement.
In the slow case, we clear it from the TLB using the \texttt{invlpg} instruction and access a virtual alias of the probe address's physical location.
This ensures that any timing differences originate only from the TLB, with the address in the L1d cache in both cases.
For each case, we take \num{16384} samples per padding level with the threads pinned to separate physical cores.
We observe a golden range of \num{39} \texttt{add}-cycles, starting at padding level \num{139}, meaning the padded \AttackName can reliably detect the timing difference.
If we additionally flush the relevant page table entries to DRAM with \texttt{clflush} before probing, we observe perfect accuracy even without padding, and a golden range of \num{178} \texttt{add}-cycles.

\subsection{\AttackName in Lua}\label{subsec:attackname-in-lua}

While previous experiments rely on direct access to binary code or simple environments like WebAssembly to implement \AttackName, the low-level capabilities these provide are not strictly required.
To demonstrate that \AttackName remains effective in a higher-level language, we devise an experiment with unpadded \AttackName for the LuaJIT environment (version 2.1.1703358377).
It follows the same structure as the WebAssembly implementation in \Cref{subsec:attackname-in-the-browser}, with the primitive implemented entirely in Lua and a simple C library to prepare the target address's cache state.
Additionally, we use LuaJIT's foreign function interface to create a shared memory region.
As with WebAssembly, we omit cache state preparation for $w$ and $s$.
We evaluate this experiment in the L1 vs.\ DRAM configuration, running on an AMD \labthirtyCPU with the threads on separate physical cores.
Overall, we take \num{1000000} samples each for the cached and the uncached case.
We observe just \num{120} misclassifications in the cached case and \num{150} in the uncached case, resulting in a combined $F_1$ score of $> 0.999$.
This shows that \AttackName is feasible in higher-level languages.



\section{Discussion}\label{sec:discussion}
In this section, we discuss \AttackName's limitations, potential mitigations, and related work.

\subsection{Limitations}\label{subsec:limitations}

\paragrabf{No Direct Timing Measurement}
Compared to a precise hardware timer, \AttackName's primary limitation is that it does not measure timing directly.
Instead, it compares two timings and reveals which is faster.
In cache attacks, this is typically not an issue.
Attacks like Flush+Reload or Prime+Probe reveal only whether an address was accessed, so absolute timing values are irrelevant.
However, this complicates attacks with more than two states to distinguish (cf. \Cref{subsec:determining-llc-slices}, where there are 12).
In this scenario, attacks using \AttackName need multiple measurements, whereas a precise timer can, in theory, distinguish them with a single sample.

\paragrabf{Precisely Controlled Padding}
The padded variant requires a mechanism for selecting a precisely controlled delay sequence.
Our native and WebAssembly implementations generate the required dependency chain at runtime.
Many higher-level languages do not expose comparable control over generated code.
Such environments can instead use loop-based padding or ship a finite collection of precompiled padding functions.
Both approaches may introduce additional jitter or provide coarser calibration steps.

\subsection{Related Work}\label{subsec:related-work}

\paragrabf{Timerless Cache Side Channels}
In environments without a timer, a standard approach for measuring timing differences is to create a timer using counting threads~\cite{Lipp2016,Schwarz2017Timers,Schwarz2020malware}.
We compare this approach to \AttackName in \Cref{subsec:comparison-to-counter-threads}.
The idea of using race conditions with thread-level parallelism for revealing memory access latencies, and thus the foundational principle of \AttackName, was proposed by Purnal \etal under the name \emph{architectural reordering}~\cite{Purnal2023showtime}.
Our work extensively evaluates and extends this idea.
Other fully timerless cache side channels typically rely on architecture-specific properties.
For instance, the abort semantics of Intel TSX may reveal cache state information~\cite{Disselkoen2017prime,kim2022dprime,Hornetz2026tdxray}.
Similarly, weak coherence between the data and instruction caches on many RISC architectures is known to enable timerless attack primitives~\cite{Guanciale2016Cache,Jin2025Ghostcache,Thomas2025Exfilstate,thomas2026riscy}.
Zhang \etal show that the \texttt{umonitor} and \texttt{umwait} instructions on x86-based architectures provide unprivileged users with architectural feedback on cache activity~\cite{Zhang2023MWAIT}.
Yu \etal create S${}^2$C, a timerless cache side channel for Apple's M1 architecture based on hardware synchronization instructions~\cite{Yu2023Synchronization}.
Thomas \etal propose ExfilState, a fuzzing-based technique for discovering architectural cache side channels on ARM~\cite{Thomas2025Exfilstate}.
In a follow-up work, Thomas~\etal\cite{thomas2026riscy} show that cache coherency effects in RISC CPUs can reveal information on the instruction cache's state.

\paragrabf{Attacks with Imprecise Timers}
In addition to fully timerless attacks like \AttackName, a large body of work exists on cache attacks with imprecise timers.
This line of work is orthogonal to \AttackName, as the limiting factor in cache attacks is often the resolution of a timer, not its full absence.
Many rely on \emph{amplification}, i.e., magnifying the timing differences between the cached and uncached states until they are measurable even with a coarse-grained timer.
A seminal work is \texttt{leaky.page}~\cite{Roettger2021SpectreJS}, which introduces an amplification primitive that relies on the Tree-PLRU cache replacement policy used in the L1 caches of many microarchitectures.
Purnal \etal extend this primitive, significantly increasing its amplification ratio~\cite{Purnal2023showtime}, and propose another amplification primitive using non-temporal prefetch instructions on x86.
Katzman \etal propose an amplification method using logic gates operating on the cache state~\cite{Katzman2023gates}.
It propagates the cache state of a target address in a tree-like structure, creating an amplified timing difference when accessing the leaves sequentially.
Morgan \etal introduce comparator gates, which create a race condition between two memory accesses during speculative execution, encoding the result in the cache~\cite{Morgan2025slice}.
This amplifies a timing difference on the order of individual cycles to a difference between L1 and DRAM\@.
Xiao \etal propose Hacky Racers~\cite{xiao2023hacky}, a primitive relying on instruction-level parallelism introduced by out-of-order execution.
It encodes the relative completion order of two parallelizable instruction paths in the cache and uses amplification to make it accessible with a coarse-grained timer.
An orthogonal approach to amplification is to average many direct measurements with a coarse-grained timer.
This requires many samples per target and was demonstrated mainly in remote attacks~\cite{Bernstein2005,Schwarz2019netspectre,Schwarzl2021dynamic}.

\subsection{Mitigations}\label{subsec:mitigations}


\paragrabf{Eliminating Shared Memory}
Shared memory, a requirement for \AttackName, can often be replaced with message passing via channels.
This alone does not fully mitigate \AttackName, as it could, in theory, also be implemented with channels by replacing writes to $s$ and $w$ with channel writes.
However, channels are often managed components, e.g., by the operating system, and thus provide a surface for implementing mitigations.
In many web browsers, for instance, message passing is already rate-limited to prevent the construction of counting threads~\cite{Kohlbrenner2016,Schwarz2017Timers}.
To further mitigate \AttackName, channels can add a random delay before delivering a message.
This makes the race-condition outcome random, preventing attackers from inferring usable information from a single measurement.
Another mitigation is to grant write access to a shared memory segment to only one thread at a time.
\AttackName requires $w$ to be writable by both threads, which such a configuration prevents.

\paragrabf{Eliminating True Parallelism}
Many programs follow an asynchronous programming model, forgoing true parallelism in favor of weaker, interleaving-based concurrency.
While this does not bring the full performance benefits of parallelism, it can unblock expensive operations, allowing programs to continue executing while waiting (e.g., for the operating system).
A prominent example is asynchronous functions in languages such as JavaScript.
Implementing \AttackName is not possible in such an environment, as it requires parallelism.

\section{Conclusion}\label{sec:conclusion}
We presented \AttackName, a timerless probing primitive for distinguishing microarchitectural latencies.
In contrast to other timerless primitives, \AttackName is portable across microarchitectures and relies only on thread-level parallelism and shared memory.
We performed an extensive evaluation of \AttackName on \PlatformCount CPUs, demonstrating that it can reliably separate microarchitectural latency differences on all of them.
Additionally, we showed that \AttackName works in a web browser, with a temporal accuracy approaching that of our native implementations.
Our case studies showed that \AttackName applies in a wide variety of contexts and can, in some instances, distinguish latency differences on the order of individual cycles.
Overall, \AttackName shows that timing-based microarchitectural attacks, even with fine-grained latency differences, do not inherently require a timer.



\FloatBarrier

\bibliographystyle{plain}
\bibliography{main}
\newpage

\appendix
\crefalias{section}{appendix}
\FloatBarrier

\section{Per-Machine Experimental Results}\label{sec:precision-analysis}
We present the per-machine experimental results from \Cref{subsec:delta-sync} and \Cref{subsec:measuring-memory-access-latencies} in \Cref{tab:precision} and \Cref{tab:precision_ht}.
$F^{fast}_{slow}$ denotes the $F_1$ score of the unpadded \AttackName in a specific setting, and $G^{fast}_{slow}$ denotes the width of the golden range with padded \AttackName in \texttt{add}-cycles.
$CTR$ is the average edge width of a counting thread, and $\Delta^{L1}_{DRAM}$ is the latency difference between an L1 hit and a DRAM access as measured in \Cref{subsec:delta-sync}.
\Cref{tab:precision_wasm} lists the results of our WASM-based implementation on x86\_64.

\section{Reliable Cache Eviction}\label{sec:reliable-cache-eviction}
Many experiments in this paper rely on eviction sets to prime the cache state, e.g., to evict an address from the L2 or L3 cache.
Unless otherwise specified, as in \Cref{subsec:timerless-evict-reload}, we construct these eviction sets using the physical addresses of both the target and eviction-set addresses.
Our eviction sets consist of $n + 2$ set-index-congruent addresses for evicting the L1 cache, and $n + 4$ addresses for evicting the L2 cache, where $n$ is the associativity of the targeted cache level.
The overfill provides additional robustness.
On CPUs where the L1 cache has higher associativity than the L2, such as Intel's Skylake architectures, we use the larger L1 associativity for $n$ when evicting the L2 cache.
To apply an eviction set, we first access the target address to cache it in the L1 cache.
Then, we access all eviction-set addresses in a linear pass, then in multiple passes in random order, and then again in a reverse-order linear pass.
To our knowledge, all L1 caches in our evaluation are directly indexed, i.e., there is no complex hash function for determining the cache set index.
Testing for congruency is therefore trivial.
With the L2 cache, we assume direct indexing for all CPUs except for Apple's M1 and M2 architectures and AMD's Zen 3, 4, and 5 architectures.
For M1, we rely on the set index function reverse-engineered by Yu et al.~\cite{Yu2023Synchronization}, and for M2, we find L2 eviction sets using the algorithm described by Vila et al.~\cite{Vila2019}.
For AMD, we use the set index functions reverse-engineered by Bagia et al.~\cite{bagia2025close} and Zhang et al.~\cite{Zhang2025Taming} to determine cache set congruency.
Where possible, we rely on cache management instructions to flush addresses from the LLC\@.
On architectures where such instructions are privileged, we use a simple kernel module to grant user-mode processes access.
We still perform LLC eviction with eviction sets only on the StarFive JH7100 and the Microchip MPFS025T.
Based on publicly available RTL code~\cite{sifive-inclusivecache}, we assume direct set indexing and construct the LLC eviction sets in the same way as L2 eviction sets on other platforms.

\begin{table*}[b]
  \caption{\AttackName characteristics with \tprobe on core 0 and \tref on the SMT sibling core. We list \rdelta{sync}, $\Delta^{L1}_{DRAM}$, and the average edge width of a counting thread (CTR). Additionally, it shows the $F_1$ score of unpadded race probing ($F^{fast}_{slow}$) and the golden range's length for padded race probing ($G^{fast}_{slow}$). All values, except the $F_1$ scores, are measured in \texttt{add}-cycles.}
  \label{tab:precision_ht}
  \centering
  \footnotesize
\renewcommand{\arraystretch}{0.75}
\begin{tabular}{cll@{\hspace{1em}}rrr@{\hspace{2.5em}}cc@{\hspace{2.5em}}cccc}
\toprule & \multirow[c]{2}{*}[-3pt]{\textbf{CPU}} & \multirow[c]{2}{*}[-3pt]{\textbf{$\mu$Architecture}} &&&& \multicolumn{2}{c}{\hspace*{-2.5em}\textbf{Unpadded}} & \multicolumn{4}{c}{\hspace*{-\tabcolsep}\textbf{Padded}} \vspace{3.5pt}\\
& &  & $\Delta_{sync}$ & $\Delta^{L1}_{DRAM}$ & CTR & $F^{L1}_{DRAM}$ & $F^{LLC}_{DRAM}$ & $G^{L1}_{DRAM}$ & $G^{L1}_{L2}$ & $G^{L1}_{L3}$ & $G^{LLC}_{DRAM}$ \vspace{1pt}\\\midrule
\multirow[c]{19}{*}{\rotatebox[origin=c]{90}{\textbf{x86}\hspace{11em}}} & AMD EPYC 7252 & Zen 2 (Rome) & $43 \pm 1$ & 329 & 74 & 1.00 & 0.99 & 41 & 5 & 31 & 12 \\\cmidrule(lr){2-12}
 & AMD Ryzen 7 5700G & Zen 3 (Cezanne) & $55 \pm 22$ & 322 & 90 & 0.99 & 0.77 & 6 & $\cdot$ & $\cdot$ & $\cdot$ \\\cmidrule(lr){2-12}
 & AMD Ryzen 9 7950X3D & Zen 4 (Raphael) & $38 \pm 29$ & 496 & 95 & 0.99 & 0.61 & 5 & $\cdot$ & $\cdot$ & $\cdot$ \\\cmidrule(lr){2-12}
 & AMD EPYC 9124 & Zen 4 (Genoa) & $50 \pm 19$ & 324 & 96 & 0.99 & 0.23 & 9 & $\cdot$ & 2 & $\cdot$ \\\cmidrule(lr){2-12}
 & AMD Ryzen 9 9950X & Zen 5 (Granite Ridge) & $76 \pm 31$ & 462 & 249 & 0.99 & 0.93 & 37 & $\cdot$ & $\cdot$ & $\cdot$ \\\cmidrule(lr){2-12}
 & Intel Core i3-2120 & Sandy Bridge & $47 \pm 8$ & 200 & 50 & 1.00 & 1.00 & 36 & $\cdot$ & $\cdot$ & 43 \\\cmidrule(lr){2-12}
 & Intel Core i3-5010U & Broadwell & $76 \pm 7$ & 156 & 43 & 0.99 & 0.99 & 18 & $\cdot$ & $\cdot$ & 14 \\\cmidrule(lr){2-12}
 & Intel Xeon E3-1505M v5 & Skylake & $35 \pm 13$ & 170 & 42 & 1.00 & 0.81 & 20 & $\cdot$ & $\cdot$ & $\cdot$ \\\cmidrule(lr){2-12}
 & Intel Pentium Gold G5400T & Coffee Lake & $44 \pm 8$ & 205 & 41 & 1.00 & 0.94 & 39 & $\cdot$ & 25 & $\cdot$ \\\cmidrule(lr){2-12}
 & Intel Core i7-8700 & Coffee Lake & $35 \pm 18$ & 272 & 25 & 1.00 & 0.24 & 39 & $\cdot$ & 39 & $\cdot$ \\\cmidrule(lr){2-12}
 & Intel Core i7-10710U & Comet Lake & $35 \pm 15$ & 305 & 43 & 1.00 & 0.92 & 20 & $\cdot$ & $\cdot$ & $\cdot$ \\\cmidrule(lr){2-12}
 & Intel Core i7-1185G7 & Willow Cove & $54 \pm 21$ & 513 & 41 & 0.99 & 0.88 & 32 & $\cdot$ & 8 & $\cdot$ \\\cmidrule(lr){2-12}
 & Intel Xeon Gold 6346 & Sunny Cove & $47 \pm 25$ & 233 & 47 & 1.00 & 1.00 & 27 & $\cdot$ & 26 & 31 \\\cmidrule(lr){2-12}
 & Intel Core i7-11700 & Cypress Cove & $53 \pm 20$ & 316 & 46 & 0.99 & 0.98 & 28 & $\cdot$ & 18 & 6 \\\cmidrule(lr){2-12}
 & Intel Core i9-12900H & Golden Cove (P) & $68 \pm 22$ & 345 & 49 & 1.00 & 0.99 & 40 & $\cdot$ & 14 & 53 \\\cmidrule(lr){2-12}
 & Intel Core i5-13420H & Raptor Cove (P) & $71 \pm 18$ & 345 & 56 & 1.00 & 0.99 & 45 & $\cdot$ & 18 & 50 \\\cmidrule(lr){2-12}
 & Intel Core i9-13900K & Raptor Cove (P) & $70 \pm 15$ & 456 & 58 & 0.99 & 0.94 & 43 & $\cdot$ & 25 & $\cdot$ \\\cmidrule(lr){2-12}
 & Intel Xeon Gold 6526Y & Raptor Cove & $54 \pm 30$ & 357 & 58 & 1.00 & 0.10 & 42 & $\cdot$ & 35 & $\cdot$ \\\cmidrule(lr){2-12}
 & Intel Xeon 6736P & Redwood Cove & $60 \pm 17$ & 681 & 50 & 1.00 & 0.94 & 45 & $\cdot$ & 39 & $\cdot$ \\\midrule
\multirow[c]{1}{*}{\textbf{LA64}} & Loongson 3A6000 & 3A6000 & $22 \pm 8$ & 183 & 3 & 1.00 & 0.00 & 14 & $\cdot$ & 14 & $\cdot$ \\
\bottomrule
\end{tabular}

\end{table*}
\begin{table*}[t]
  \caption{\AttackName characteristics with \tprobe on core 0 and \tref on the physical core with the next-highest ID. We list \rdelta{sync}, $\Delta^{L1}_{DRAM}$, and the average edge width of a counting thread (CTR). Additionally, it shows the $F_1$ score of unpadded race probing ($F^{fast}_{slow}$) and the golden range's length for padded race probing ($G^{fast}_{slow}$). All values, except the $F_1$ scores, are measured in \texttt{add}-cycles.}
  \label{tab:precision}
  \centering
  \footnotesize
\renewcommand{\arraystretch}{0.75}
\begin{tabular}{cll@{\hspace{1em}}rrr@{\hspace{2.5em}}cc@{\hspace{2.5em}}cccc}
\toprule & \multirow[c]{2}{*}[-3pt]{\textbf{CPU}} & \multirow[c]{2}{*}[-3pt]{\textbf{$\mu$Architecture}} &&&& \multicolumn{2}{c}{\hspace*{-2.5em}\textbf{Unpadded}} & \multicolumn{4}{c}{\hspace*{-\tabcolsep}\textbf{Padded}} \vspace{3.5pt}\\
& &  & $\Delta_{sync}$ & $\Delta^{L1}_{DRAM}$ & CTR & $F^{L1}_{DRAM}$ & $F^{LLC}_{DRAM}$ & $G^{L1}_{DRAM}$ & $G^{L1}_{L2}$ & $G^{L1}_{L3}$ & $G^{LLC}_{DRAM}$ \vspace{1pt}\\\midrule
\multirow[c]{8}{*}{\rotatebox[origin=c]{90}{\textbf{ARM}\hspace{4em}}} & \multirow[c]{2}{*}[-2.25pt]{Apple M1} & Icestorm (E) & $122 \pm 1$ & 236 & 86 & 0.99 & 0.99 & 177 & 8 & 88 & 87 \\\cmidrule(lr){3-12}
 &  & Firestorm (P) & $142 \pm 1$ & 342 & 105 & 1.00 & 0.99 & 230 & 53 & 122 & 44 \\\cmidrule(lr){2-12}
 & \multirow[c]{2}{*}[-2.25pt]{Apple M2} & Blizzard (E) & $123 \pm 6$ & 290 & 97 & 0.97 & 0.99 & 123 & $\cdot$ & 24 & 65 \\\cmidrule(lr){3-12}
 &  & Avalanche (P) & $142 \pm 7$ & 371 & 112 & 0.97 & 0.99 & 92 & $\cdot$ & $\cdot$ & 355 \\\cmidrule(lr){2-12}
 & Broadcom BCM2837 & Cortex-A53 & $49 \pm 0$ & 161 & 35 & 0.99 & 0.99 & 50 & 13 & -- & 37 \\\cmidrule(lr){2-12}
 & Broadcom BCM2712 & Cortex-A76 & $143 \pm 61$ & 322 & 172 & 0.99 & 0.99 & 161 & $\cdot$ & $\cdot$ & 234 \\\cmidrule(lr){2-12}
 & Huawei Kunpeng 920 & TaiShan-v110 & $243 \pm 1$ & 271 & 91 & 0.75 & 0.75 & 104 & $\cdot$ & $\cdot$ & 92 \\\cmidrule(lr){2-12}
 & Ampere Altra Q64-30 & Altra & $293 \pm 20$ & 251 & 242 & 0.67 & 0.67 & 99 & $\cdot$ & $\cdot$ & 120 \\\midrule
\multirow[c]{31}{*}{\rotatebox[origin=c]{90}{\textbf{x86}\hspace{18em}}} & AMD EPYC 7252 & Zen 2 (Rome) & $78 \pm 0$ & 329 & 205 & 0.99 & 0.99 & 78 & 10 & 34 & 44 \\\cmidrule(lr){2-12}
 & AMD Ryzen 7 5700G & Zen 3 (Cezanne) & $84 \pm 0$ & 322 & 162 & 0.99 & 0.99 & 71 & $\cdot$ & 33 & 37 \\\cmidrule(lr){2-12}
 & AMD Ryzen 9 7950X3D & Zen 4 (Raphael) & $87 \pm 3$ & 496 & 183 & 1.00 & 0.99 & 66 & $\cdot$ & 6 & 38 \\\cmidrule(lr){2-12}
 & AMD EPYC 9124 & Zen 4 (Genoa) & $85 \pm 5$ & 324 & 195 & 0.99 & 0.84 & 49 & $\cdot$ & $\cdot$ & $\cdot$ \\\cmidrule(lr){2-12}
 & AMD Ryzen 9 9950X & Zen 5 (Granite Ridge) & $140 \pm 1$ & 462 & 201 & 0.99 & 0.99 & 125 & 10 & 47 & 78 \\\cmidrule(lr){2-12}
 & Intel Core i3-2120 & Sandy Bridge & $71 \pm 8$ & 200 & 102 & 0.99 & 0.99 & 61 & 2 & 7 & 62 \\\cmidrule(lr){2-12}
 & Intel Core i5-4690 & Haswell & $95 \pm 1$ & 281 & 116 & 0.99 & 1.00 & 89 & $\cdot$ & 13 & 100 \\\cmidrule(lr){2-12}
 & Intel Core i3-5010U & Broadwell & $88 \pm 3$ & 156 & 117 & 0.99 & 0.99 & 88 & 2 & 12 & 54 \\\cmidrule(lr){2-12}
 & Intel Xeon E3-1505M v5 & Skylake & $90 \pm 1$ & 170 & 126 & 1.00 & 1.00 & 97 & 2 & 21 & 68 \\\cmidrule(lr){2-12}
 & Intel Core i5-6400T & Skylake & $95 \pm 5$ & 88 & 134 & 0.68 & 0.68 & 80 & $\cdot$ & 19 & 50 \\\cmidrule(lr){2-12}
 & Intel Pentium G4400T & Skylake & $88 \pm 4$ & 190 & 120 & 0.99 & 0.99 & 92 & $\cdot$ & $\cdot$ & 63 \\\cmidrule(lr){2-12}
 & Intel Pentium Gold G5400T & Coffee Lake & $92 \pm 3$ & 205 & 146 & 1.00 & 1.00 & 91 & $\cdot$ & 22 & 59 \\\cmidrule(lr){2-12}
 & Intel Core i7-8700 & Coffee Lake & $96 \pm 3$ & 272 & 114 & 1.00 & 0.99 & 95 & $\cdot$ & 27 & 54 \\\cmidrule(lr){2-12}
 & Intel Core i7-9700K & Coffee Lake & $91 \pm 5$ & 298 & 156 & 1.00 & 0.99 & 103 & $\cdot$ & 22 & 73 \\\cmidrule(lr){2-12}
 & Intel Core i7-10710U & Comet Lake & $95 \pm 3$ & 305 & 224 & 0.99 & 0.99 & 90 & $\cdot$ & 21 & 57 \\\cmidrule(lr){2-12}
 & Intel Core i7-1185G7 & Willow Cove & $126 \pm 17$ & 513 & 231 & 0.99 & 0.99 & 103 & $\cdot$ & 9 & 55 \\\cmidrule(lr){2-12}
 & Intel Xeon Gold 6346 & Sunny Cove & $151 \pm 6$ & 233 & 300 & 1.00 & 0.69 & 169 & 12 & 50 & 143 \\\cmidrule(lr){2-12}
 & Intel Core i7-11700 & Cypress Cove & $134 \pm 12$ & 316 & 241 & 0.99 & 1.00 & 131 & $\cdot$ & $\cdot$ & 83 \\\cmidrule(lr){2-12}
 & \multirow[c]{2}{*}[-2.25pt]{Intel Core i9-12900H} & Gracemont (E) & $147 \pm 3$ & 295 & 194 & 1.00 & 1.00 & 147 & 14 & 68 & 74 \\\cmidrule(lr){3-12}
 &  & Golden Cove (P) & $145 \pm 22$ & 345 & 244 & 0.99 & 0.99 & 133 & $\cdot$ & 24 & 82 \\\cmidrule(lr){2-12}
 & \multirow[c]{2}{*}[-2.25pt]{Intel Core i5-13420H} & Gracemont (E) & $156 \pm 4$ & 288 & 198 & 0.99 & 0.99 & 157 & 7 & 33 & 113 \\\cmidrule(lr){3-12}
 &  & Raptor Cove (P) & $142 \pm 15$ & 345 & 258 & 0.99 & 0.99 & 92 & $\cdot$ & 30 & 21 \\\cmidrule(lr){2-12}
 & \multirow[c]{2}{*}[-2.25pt]{Intel Core i9-13900K} & Gracemont (E) & $170 \pm 10$ & 392 & 166 & 1.00 & 1.00 & 167 & 14 & 52 & 97 \\\cmidrule(lr){3-12}
 &  & Raptor Cove (P) & $179 \pm 13$ & 456 & 357 & 0.99 & 0.99 & 156 & $\cdot$ & 31 & 61 \\\cmidrule(lr){2-12}
 & Intel Xeon Gold 6526Y & Raptor Cove & $143 \pm 8$ & 357 & 355 & 0.99 & 0.99 & 169 & 8 & 65 & 93 \\\cmidrule(lr){2-12}
 & Intel Core Ultra 9 288V & Lion Cove & $153 \pm 13$ & 279 & 286 & 0.99 & 0.99 & 135 & $\cdot$ & 38 & 71 \\\cmidrule(lr){2-12}
 & \multirow[c]{2}{*}[-2.25pt]{Intel Core Ultra 9 285K} & Skymont (E) & $98 \pm 0$ & 475 & 86 & 1.00 & 1.00 & 97 & $\cdot$ & 10 & 33 \\\cmidrule(lr){3-12}
 &  & Lion Cove (P) & $162 \pm 17$ & 577 & 266 & 0.99 & 0.99 & 234 & $\cdot$ & 40 & 172 \\\cmidrule(lr){2-12}
 & Intel Xeon 6736P & Redwood Cove & $403 \pm 23$ & 681 & 792 & 0.99 & 0.99 & 396 & $\cdot$ & 138 & 193 \\\cmidrule(lr){2-12}
 & \multirow[c]{2}{*}[-2.25pt]{Intel Core Ultra 5 325} & Darkmont (E) & $90 \pm 7$ & 121 & 86 & 0.94 & 0.93 & 63 & $\cdot$ & $\cdot$ & 47 \\\cmidrule(lr){3-12}
 &  & Cougar Cove (P) & $151 \pm 25$ & 284 & 202 & 1.00 & 0.99 & 150 & $\cdot$ & 28 & 88 \\\midrule
\multirow[c]{4}{*}{\rotatebox[origin=c]{90}{\textbf{RISC-V}\hspace{1em}}} & StarFive JH7100 & U74 & $27 \pm 15$ & 129 & 55 & 0.98 & 0.59 & 2 & $\cdot$ & -- & $\cdot$ \\\cmidrule(lr){2-12}
 & Microchip MPFS025T & U54 & $18 \pm 0$ & 75 & 54 & 0.99 & 0.66 & 19 & $\cdot$ & -- & $\cdot$ \\\cmidrule(lr){2-12}
 & T-Head TH1520 & C910 & $56 \pm 0$ & 239 & 86 & 0.99 & 0.99 & 48 & 29 & -- & 48 \\\cmidrule(lr){2-12}
 & SpacemiT K1 & X60 & $75 \pm 7$ & 145 & 216 & 0.99 & 1.00 & 55 & 2 & -- & 45 \\\midrule
\multirow[c]{2}{*}{\rotatebox[origin=c]{90}{\hspace{.5em}\textbf{LA64}}} & Loongson 3A5000 & 3A5000 & $186 \pm 2$ & 309 & 123 & 0.99 & 1.00 & 150 & $\cdot$ & 5 & 140 \\\cmidrule(lr){2-12}
 & Loongson 3A6000 & 3A6000 & $120 \pm 2$ & 183 & 100 & 1.00 & 1.00 & 119 & $\cdot$ & 33 & 84 \\
\bottomrule
\end{tabular}

\end{table*}
\begin{table*}[t]
    \caption{\AttackName performance in WebAssembly (Chromium v147.0) in the L1 vs. DRAM setting on x86\_64. \rdelta{sync} and the width of the golden range $G^{L1}_{DRAM}$ are measured in \texttt{add}-cycles.}
    \label{tab:precision_wasm}
    \centering
    \footnotesize
\renewcommand{\arraystretch}{0.75}
\begin{tabular}{ll@{\hspace{2em}}rr@{\hspace{3em}}cc@{\hspace{3em}}cc}
\toprule \multirow[c]{2}{*}[-3pt]{\textbf{CPU}} & \multirow[c]{2}{*}[-3pt]{\textbf{$\mu$Architecture}} & \multicolumn{2}{c}{\hspace*{-3em}\makebox[0pt]{$\Delta_{sync}$}} & \multicolumn{2}{c}{\hspace*{-3em}\makebox[0pt]{$F^{L1}_{DRAM}$ (unpadded)}} & \multicolumn{2}{c}{\hspace*{-\tabcolsep}\makebox[0pt]{$G^{L1}_{DRAM}$ (padded)}} \vspace{1pt}\\\cmidrule(lr{\dimexpr 3em+0.5em\relax}){3-4}\cmidrule(lr{\dimexpr 3em+0.5em\relax}){5-6}\cmidrule(lr{\dimexpr \tabcolsep+0.5em\relax}){7-8}
 & & SMT & Cross-core & SMT & Cross-core & SMT & Cross-core \vspace{1pt}\\\midrule
AMD EPYC 7252 & Zen 2 (Rome) & $50 \pm 4$ & $101 \pm 33$ & 0.99 & 0.99 & 38 & 67 \\\cmidrule(lr){1-8}
AMD Ryzen 7 5700G & Zen 3 (Cezanne) & $48 \pm 20$ & $126 \pm 3$ & 0.99 & 0.99 & 44 & 113 \\\cmidrule(lr){1-8}
AMD Ryzen 9 7950X3D & Zen 4 (Raphael) & $49 \pm 21$ & $88 \pm 13$ & 0.99 & 0.99 & 50 & 69 \\\cmidrule(lr){1-8}
AMD EPYC 9124 & Zen 4 (Genoa) & $48 \pm 18$ & $120 \pm 1$ & 0.99 & 0.99 & 42 & 128 \\\cmidrule(lr){1-8}
AMD Ryzen 9 9950X & Zen 5 (Granite Ridge) & $91 \pm 7$ & $165 \pm 7$ & 0.99 & 0.99 & 69 & 137 \\\cmidrule(lr){1-8}
Intel Core i3-2120 & Sandy Bridge & $24 \pm 12$ & $84 \pm 30$ & 0.99 & 0.98 & 22 & 36 \\\cmidrule(lr){1-8}
Intel Core i5-4690 & Haswell & -- & $83 \pm 1$ & -- & 0.99 & -- & 69 \\\cmidrule(lr){1-8}
Intel Core i3-5010U & Broadwell & $28 \pm 14$ & $108 \pm 4$ & 0.99 & 0.99 & 25 & 61 \\\cmidrule(lr){1-8}
Intel Xeon E3-1505M v5 & Skylake & $33 \pm 20$ & $89 \pm 7$ & 0.99 & 0.99 & 41 & 66 \\\cmidrule(lr){1-8}
Intel Core i5-6400T & Skylake & -- & $99 \pm 5$ & -- & 0.97 & -- & 50 \\\cmidrule(lr){1-8}
Intel Pentium G4400T & Skylake & -- & $91 \pm 43$ & -- & 0.99 & -- & 26 \\\cmidrule(lr){1-8}
Intel Pentium Gold G5400T & Coffee Lake & $36 \pm 17$ & $128 \pm 6$ & 0.99 & 0.99 & 30 & 98 \\\cmidrule(lr){1-8}
Intel Core i7-8700 & Coffee Lake & $36 \pm 17$ & $110 \pm 26$ & 0.99 & 0.99 & 28 & 75 \\\cmidrule(lr){1-8}
Intel Core i7-9700K & Coffee Lake & -- & $106 \pm 4$ & -- & 0.98 & -- & 82 \\\cmidrule(lr){1-8}
Intel Core i7-10710U & Comet Lake & $32 \pm 19$ & $120 \pm 12$ & 0.99 & 0.99 & 11 & 29 \\\cmidrule(lr){1-8}
Intel Core i7-1185G7 & Willow Cove & $46 \pm 11$ & $122 \pm 18$ & 0.99 & 0.99 & 41 & 90 \\\cmidrule(lr){1-8}
Intel Xeon Gold 6346 & Sunny Cove & $36 \pm 17$ & $327 \pm 113$ & 0.99 & 0.99 & 37 & 316 \\\cmidrule(lr){1-8}
Intel Core i7-11700 & Cypress Cove & $39 \pm 6$ & $130 \pm 5$ & 0.99 & 0.99 & 11 & 88 \\\cmidrule(lr){1-8}
\multirow[c]{2}{*}[-2.25pt]{Intel Core i9-12900H} & Gracemont (E) & -- & $140 \pm 5$ & -- & 0.99 & -- & 125 \\\cmidrule(lr){2-8}
 & Golden Cove (P) & $55 \pm 16$ & $229 \pm 37$ & 0.99 & 0.99 & 27 & 139 \\\cmidrule(lr){1-8}
\multirow[c]{2}{*}[-2.25pt]{Intel Core i5-13420H} & Gracemont (E) & -- & $128 \pm 6$ & -- & 0.99 & -- & 73 \\\cmidrule(lr){2-8}
 & Raptor Cove (P) & $55 \pm 16$ & $193 \pm 17$ & 0.99 & 0.99 & 17 & 155 \\\cmidrule(lr){1-8}
\multirow[c]{2}{*}[-2.25pt]{Intel Core i9-13900K} & Gracemont (E) & -- & $168 \pm 4$ & -- & 0.99 & -- & 123 \\\cmidrule(lr){2-8}
 & Raptor Cove (P) & $54 \pm 15$ & $276 \pm 16$ & 0.99 & 0.99 & 31 & 188 \\\cmidrule(lr){1-8}
Intel Xeon Gold 6526Y & Raptor Cove & $49 \pm 28$ & $303 \pm 21$ & 0.99 & 0.99 & 33 & 191 \\\cmidrule(lr){1-8}
Intel Core Ultra 9 288V & Lion Cove & -- & $173 \pm 23$ & -- & 0.99 & -- & 152 \\\cmidrule(lr){1-8}
\multirow[c]{2}{*}[-2.25pt]{Intel Core Ultra 9 285K} & Skymont (E) & -- & $180 \pm 2$ & -- & 0.97 & -- & 142 \\\cmidrule(lr){2-8}
 & Lion Cove (P) & -- & $326 \pm 31$ & -- & 0.99 & -- & 218 \\\cmidrule(lr){1-8}
Intel Xeon 6736P & Redwood Cove & $53 \pm 20$ & $974 \pm 31$ & 0.99 & 0.99 & 35 & 361 \\\cmidrule(lr){1-8}
\multirow[c]{2}{*}[-2.25pt]{Intel Core Ultra 5 325} & Darkmont (E) & -- & $130 \pm 100$ & -- & 0.99 & -- & 14 \\\cmidrule(lr){2-8}
 & Cougar Cove (P) & -- & $172 \pm 12$ & -- & 0.99 & -- & 128 \\
\bottomrule
\end{tabular}

\end{table*}

\end{document}
